\chapter{PRAGMATIQUE}
\label{chapter-12}
\markright{Pragmatique}
\origpage{350}{336}
\begin{keybox}[Plan du chapitre]
\small\setstretch{1.01}
\setlist{topsep=.12em,itemsep=.12em,parsep=0pt}
\renewcommand{\thempfootnote}{校\arabic{editorialnote}}
\textbf{I.\quad LA LINGUISTIQUE DE L’ÉNONCIATION}
\begin{enumerate}[label=\arabic*.]
\item Énonciation et situation d’énonciation
\begin{enumerate}[label=\textcircled{\scriptsize\arabic*}]
\item Énonciation, énoncé, situation d’énonciation
\item Énoncé coupé de la situation d’énonciation
\item Énoncé ancré dans la situation d’énonciation
\end{enumerate}
\item Actants et circonstants de l’énonciation
\begin{enumerate}[label=\textcircled{\scriptsize\arabic*}]
\item Énonciateur, sujet parlant et locuteur\par 1. Mikhaïl Bakhtine : le dialogisme\par 2. Oswald Ducrot : la polyphonie
\item Le destinataire
\item Les circonstants
\end{enumerate}
\end{enumerate}
\textbf{II.\quad EMBRAYEURS ET DÉICTIQUES}
\begin{enumerate}[label=\arabic*.]
\item Les embrayeurs
\begin{enumerate}[label=\textcircled{\scriptsize\arabic*}]
\item Les embrayeurs subjectifs\par 1. Embrayeurs de la première personne\par 2. Embrayeurs de la deuxième personne\par 3. Modalisateurs du discours
\item Les embrayeurs spatio-temporels\par 1. Embrayeurs spatiaux\par 2. Embrayeurs temporels
\end{enumerate}
\item Les déictiques
\begin{enumerate}[label=\textcircled{\scriptsize\arabic*}]
\item Définition
\item Les formes de déixis\par 1. La déixis de personne\par 2. La déixis spatiale\par 3. La déixis temporelle\par 4. Exemple
\end{enumerate}
\item Endophore et exophore
\begin{enumerate}[label=\textcircled{\scriptsize\arabic*}]
\item L’endophore\par 1. L’anaphore
\tcbbreak
\origpage{351}{337}
\hspace*{1em}1) Définition\par\hspace*{1em}2) Fonctions\par\hspace*{1em}3) Catégories\par 2. La cataphore
\item L’exophore\par 1. Exophore mémorielle\par 2. Exophore a-mémorielle
\item Hétérogénéité des marqueurs référentiels
\end{enumerate}
\end{enumerate}
\textbf{III.\quad LA PRAGMATIQUE : DIRE, C’EST FAIRE}
\begin{enumerate}[label=\arabic*.]
\item Les énoncés constatifs et performatifs
\begin{enumerate}[label=\textcircled{\scriptsize\arabic*}]
\item La notion de performativité
\item Les expressions de la performativité\par 1. Les énoncés à performativité lexicalement dénommée\par 2. Les énoncés à performativité indiquée autrement
\item Limites de la distinction entre performatif et constatif
\end{enumerate}
\item Les actes de langage
\begin{enumerate}[label=\textcircled{\scriptsize\arabic*}]
\item La classification des \textit{speech acts}\par 1. Les actes locutoires\par 2. Les actes illocutoires\par 3. Les actes perlocutoires\par 4. Différence entre actes illocutoires et perlocutoires
\item La performativité selon Derrida
\item La classification des actes illocutoires
\item La classification des verbes performatifs
\end{enumerate}
\item La pragmatique conversationnelle
\begin{enumerate}[label=\textcircled{\scriptsize\arabic*}]
\item Oswald Ducrot : les lois du discours
\item Paul Grice : les maximes et implicature conversationnelles\ednote{章首 implicature 用了单数，但与 les、conversationnelles 不一致，应为 implicatures conversationnelles；正文对应标题亦见。}
\par 1. Le principe de coopération\par 2. Les maximes conversationnelles\par 3. Les violations volontaires\par 4. La notion d’implicature\par\hspace*{1em}1) L’implicature conventionnelle\par\hspace*{1em}2) L’implicature conversationnelle
\item Geoffrey Leech : la pragmatique de la politesse\par 1. Les maximes de politesse\par\hspace*{1em}1) Maxime de tact\par\hspace*{1em}2) Maxime de générosité\par\hspace*{1em}3) Maxime d’approbation
\tcbbreak
\origpage{352}{338}
\hspace*{1em}4) Maxime de modestie\par\hspace*{1em}5) Maxime d’accord\par\hspace*{1em}6) Maxime de sympathie\par 2. Politesse positive et politesse négative\par 3. La politesse paradoxale
\item Brown et Levinson : les \textit{Face Threatening Acts}\par 1. Face positive et face négative\par 2. Les menaces faites à la face\par 3. Les menaces faites au territoire\par 4. Les notions de \textit{face want} et \textit{face work}
\item Kerbrat-Orecchioni : les \textit{Face Flattering Acts}\par 1. FFA et anti FTA\par 2. Les adoucisseurs de FTA
\end{enumerate}
\end{enumerate}
\textbf{Exercices}\par\textbf{Pour aller plus loin}\par\textbf{Glossaire}
\end{keybox}

\origpage{353}{339}
\begin{keybox}[Introduction]
\renewcommand{\thempfootnote}{校\arabic{editorialnote}}
S’il est un point commun à toutes les langues, c’est bien d’être parlées.\ednote{此处将所有语言一概称为“被说出”的语言，需限定：自然手语并不使用口语的声音通道。若 parlées 仅宽泛指语言被使用，须与本段随后讨论的 parole 及书面语对立区别。参见本书第七章及前文关于手语的讨论。} Même s’il est vrai que dans l’antiquité gréco-latine, la rhétorique a fourni une abondante réflexion sur certains usages de la parole dans une société où l’art de l’orateur était au cœur de la vie de la cité, dans les cultures à forte tradition littéraire, la parole a longtemps été occultée par l’écrit.

Quant à la mise à l’écart de la parole par les linguistes, est-il nécessaire de rappeler que le responsable est Ferdinand de Saussure, qui, opposant la langue à la parole, a choisi la première comme objet de la linguistique moderne et, ce faisant, a établi la légitimité d’une linguistique de la langue par opposition aux « linguistiques de la parole ». Pour lui, rappelons que l’étude des phénomènes linguistiques externes ne touche pas aux éléments internes puisque « la langue est un système qui ne connaît que son ordre propre ». Tout linguiste, pour Saussure, doit faire un choix :

\begin{quote}
C’est la première bifurcation qu’on rencontre dès qu’on cherche à faire la théorie du langage. Il faut choisir entre deux routes qu’il est impossible de prendre en même temps ; elles doivent être suivies séparément. On peut à la rigueur conserver le nom de linguistique à chacune de ces deux disciplines et parler d’une linguistique de la parole. Mais il ne faudra pas la confondre avec la linguistique proprement dite, celle dont la langue est l’unique objet.
\end{quote}

En Europe, le linguiste français Émile Benveniste a décidé d’« aller au-delà du point où Saussure s’est arrêté » et a choisi une autre route en initiant une théorie de la parole et de l’énonciation. Ces théories de l’énonciation ont ensuite été prolongées par la linguistique pragmatique, d’inspiration anglo-saxonne, et dont les travaux (théorie des actes de langages, performativité, linguistique conversationnelle, linguistique de la politesse, etc.) seront présentés dans la deuxième partie dans ce chapitre.\ednote{原文 dans la deuxième partie dans ce chapitre 的后一 dans 应为 de；按本章实际编号，言语行为与会话语用学位于第三节（III），并非第二节。原文 actes de langages 通常亦写作 actes de langage。}
\end{keybox}

\section{LA LINGUISTIQUE DE L’ÉNONCIATION}
Élève d’\textbf{Antoine Meillet} à l’École Pratique des Hautes Études, Émile Benveniste (1907-1976)\ednote{1907 应为 1902；下一页同书例句亦给出 1902。此项依据本书内部对照。\ecite{56}} a enseigné lui-même dans cet établissement, ainsi qu’au Collège de France, où il a occupé la chaire de grammaire comparée de 1937 à 1969.\ednote{任教终年待核：应区分实际停止授课与正式教席任期的终点。正式任职记录仍待核。} Spécialiste de l’iranien, il s’est illustré dans la grammaire comparée des langues indo-européennes et en linguistique générale. On considère que ce sont ses \textit{Problèmes de linguistique générale} (1966 et 1974), qui, les premiers, ont introduit en France la linguistique de l’énonciation.

\subsection{Énonciation et situation d’énonciation}
Dans toute communication, aussi bien orale qu’écrite, on trouve à la fois un énoncé et une énonciation. Pour Benveniste : « C’est l’énonciation qui fait l’énoncé. »

\origpage{354}{340}
\minititle{\textcircled{\scriptsize 1} Énonciation, énoncé, situation d’énonciation}
L’énonciation est définie comme l’acte individuel de production d’un énoncé, adressé à un destinataire, dans certaines circonstances. L’énoncé, quant à lui, est le résultat linguistique, c’est-à-dire la parole prononcée ou le texte écrit. On pourrait dire que l’énoncé est le « dit » et que l’énonciation est le « dire ».

L’énoncé, saisissable par l’ouïe dans le cas de l’oral et par la vue dans celui de l’écrit, est reproductible : cela le différencie de l’énonciation, qui, en tant qu’acte individuel et unique, ne peut, par nature, être reproduite.

La situation d’énonciation, quant à elle, désigne la situation dans laquelle a été produit un énoncé. Elle permet de déterminer qui parle à qui et dans quelles circonstances. Pour reprendre la terminologie de Tesnière, l’acte d’énonciation met en scène des actants et des circonstants. Or, selon que les actants et les circonstants de la situation d’énonciation sont ou non présents dans un énoncé donné, celui-ci sera « ancré » ou bien « coupé » de la situation d’énonciation.

\minititle{\textcircled{\scriptsize 2} Énoncé coupé de la situation d’énonciation}
Un énoncé coupé de la situation d’énonciation (« non embrayé ») ne comporte aucun indice (embrayeur) permettant de repérer celle-ci. Ce sont principalement des récits au passé et à la troisième personne : une distance est établie entre ce qui est raconté et le présent de celui qui raconte, comme dans l’exemple suivant :

\textbf{Ex :} Né à Alep en 1902, naturalisé français en 1924, Benveniste est mort en 1976.

\minititle{\textcircled{\scriptsize 3} Énoncé ancré dans la situation d’énonciation}
À l’inverse, un énoncé ancré dans la situation d’énonciation (« embrayé ») comporte au moins un indice permettant de repérer celle-ci.

\textbf{Ex :} La semaine dernière, dans la salle de classe, le professeur a dit à ses étudiants : « Demain, je ne pourrai pas venir ici vous faire cours. »

L’énoncé contient cette fois-ci un certain nombre d’embrayeurs permettant de le mettre en relation avec sa propre situation d’énonciation. Ce sont :
\begin{itemize}
\item L’adverbe « demain », embrayeur temporel signifiant « dans une journée après aujourd’hui ».
\item Le pronom personnel « je », un embrayeur de la première personne désignant l’énonciateur, soit « le professeur ».
\item Le verbe « pourrai », ou plus précisément, la flexion « ai » : marque de la première personne du singulier, est également un embrayeur de la première personne désignant l’énonciateur.
\item Le pronom personnel « vous » est un embrayeur de la deuxième personne du pluriel renvoyant aux destinataires, soit « les étudiants ».
\item L’adverbe « ici », embrayeur spatial, signifiant « dans la salle de classe ».\ednote{例中含两个话语层次：“La semaine dernière”相对于转述者说话时间定位；引号内 demain、je、ici、vous 则相对于教师说话时的语境定位。不能把全部指示语统一回指到转述时刻；本页列表只分析了直接引语层。}
\end{itemize}

\origpage{355}{341}
Ce second énoncé est donc ancré dans la situation d’énonciation.

\subsection{Actants et circonstants de l’énonciation}
Comme leur nom l’indique, les circonstants correspondent aux circonstances de l’énonciation. Les actants de l’énonciation sont d’une part l’énonciateur (celui qui parle ou qui écrit), d’autre part le destinataire, c’est-à-dire celui à qui s’adresse l’énoncé, parlé ou écrit. Les « indices » permettant de repérer la participation de l’énonciateur à la situation d’énonciation, la présence du destinataire, ainsi que les circonstances de lieu et de temps dans lesquelles est produit l’énoncé, sont appelés, nous l’avons vu, « embrayeurs ».

\minititle{\textcircled{\scriptsize 1} Énonciateur, sujet parlant et locuteur}
L’énonciateur est l’actant qui dit « je ». Parfois appelé sujet de l’énonciation, il ne faut pas le confondre avec le sujet de l’énoncé :

\textbf{Ex :} Nietzsche a dit : « Dieu est mort. »

Dans l’exemple ci-dessus, le sujet de l’énoncé est « Dieu ». Mais le sujet de l’énonciation (l’énonciateur) est la personne qui produit cet énoncé, c’est-à-dire Nietzsche. L’énonciateur est appelé locuteur à l’oral, et scripteur à l’écrit. Dans le cas du discours indirect ou rapporté, les choses se compliquent encore un peu, au point que l’on peut se demander qui parle et si l’énonciateur est unique :

\textbf{Ex :} Paul m’a dit que le cours de linguistique de demain est annulé.

En linguistique, ce phénomène qui consiste à entendre dans un énoncé plusieurs « voix » distinctes est appelé dialogisme ou polyphonie.

\minititle{1. Mikhaïl Bakhtine : le dialogisme}
Le dialogisme, souvent associé à la polyphonie, est un concept développé par le philosophe et théoricien de la littérature \textbf{Mikhaïl Bakhtine} (1895-1975) dans son ouvrage \textit{Problème de la poétique de Dostoïevski} (1929),\ednote{版本待核：此处将题名“诗学”与 1929 年连用，需区分早期版本与修订本，不宜仅凭后出译本题名确定初版信息。法文书名的单复数也待相应题名页核对。两个俄文版本的题名页仍待核。\ecite{57}} et dans lequel il définit le dialogisme comme l’interaction qui se constitue entre le discours du narrateur principal et les discours d’autres personnages ou entre deux discours internes d’un personnage. Dans la partie « Du discours romanesque » de son ouvrage posthume \textit{Esthétique et théorie du roman} (1978), il précise :

\begin{lecture}{}
Il suffit d’écouter et de méditer les paroles qu’on entend partout pour affirmer ceci : dans le parler courant de tout homme vivant en société, la moitié au moins des paroles qu’il prononce sont celles d’autrui.
\end{lecture}

Le dialogisme est un procédé littéraire par lequel l’auteur peut laisser entendre une voix et une pensée indépendantes de la sienne tout en gardant une position neutre et sans qu’aucun point de vue ne soit privilégié.

\origpage{356}{342}
\minititle{2. Oswald Ducrot : la polyphonie}
Dans le chapitre 4, nous avons vu que le linguiste français \textbf{Oswald Ducrot} s’est beaucoup intéressé à la question de ce que l’on dit lorsque l’on parle. Rappelons ici les titres de certains de ses ouvrages : \textit{La preuve et le dire, Dire et ne pas dire, Le dire et le dit}. En adaptant à l’analyse conversationnelle le phénomène constaté en analyse littéraire par Mikhaïl Bakhtine, Ducrot considère qu’il n’y a pas une voix unique dans les énoncés, mais plusieurs. Citer un proverbe, faire des citations, parler au nom des autres, tout cela, pour Ducrot, relève de la polyphonie :

\textbf{Ex :} Le professeur dit : « Pour Saussure, dans la langue il n’a que des différences. »\ednote{引语中的 il n’a que 应为 il n’y a que；原文如此。}

\hspace*{1.5em}« Qui vole un oeuf, vole un boeuf. »

Pour Ducrot, le locuteur n’est pas homogène : il est un ensemble hétérogène d’emprunts, « un orchestre de voix ». Partant de la distinction effectuée par \textbf{Charles Bally}, qui, dès 1932, distinguait « sujet modal » et « sujet parlant »\origfoot{1}{Contrairement au sujet parlant qui est un être empirique, le sujet modal est une image intérieure au sens. Ce dernier n’est pas celui qui parle et pense effectivement, mais celui que le langage présente « comme celui » qui parle et pense ce qui est dit.}, Oswald Ducrot a distingué à son tour le sujet parlant (producteur empirique de l’énoncé, équivalent de l’auteur) du locuteur (instance qui prend la responsabilité de l’acte de langage, équivalent du narrateur).

\minititle{\textcircled{\scriptsize 2} Le destinataire}
Vous l’avez sûrement deviné : le destinataire est l’actant à qui l’énonciateur d’adresse.\ednote{d’adresse 应为 s’adresse，系代词误植。} Il est parfois appelé « récepteur » ou « co-énonciateur ». À l’oral, le destinataire est appelé interlocuteur, ou allocutaire. À l’écrit, on parle de lecteur. Contrairement à l’énonciateur, le destinataire peut être multiple (quand on s’adresse à plusieurs personnes à la fois).\ednote{“与发话者不同”这一限制过强：多人合署的声明等也可有共同发话主体。此处可作为单一发话者的分析假设，不宜读成对一切交际的排他规定。}

\minititle{\textcircled{\scriptsize 3} Les circonstants}
Les circonstants renvoient pour l’essentiel, aux circonstances de lieu et de temps. Un circonstant de lieu est déterminé par rapport au lieu de l’énonciation : « ici », c’est-à-dire, l’endroit où le « je » parle).\ednote{parle 后有一个无对应左括号的右括号；照录并注明排印疑误。} Un circonstant de temps est déterminé par rapport au temps de l’énonciation : « maintenant », c’est-à-dire, le moment où le « je » parle.

Nous avons déjà rapidement évoqué les embrayeurs, ces indices qui permettent d’ancrer l’énonciation. À ce terme d’embrayeur est souvent associé un autre terme, celui de « déictiques ».

\section{EMBRAYEURS ET DÉICTIQUES}
\textbf{Roman Jakobson}, dont nous avons étudié le schéma de la communication, a souligné l’importance des embrayeurs (qu’il appelle \textit{shifters}). Éléments dont le fonctionnement sémantique est inséparable de la situation d’énonciation, on les retrouve dans pratiquement tous les systèmes linguistiques, au point qu’on peut les considérer comme des « universaux » du langage.

\origpage{357}{343}
\subsection{Les embrayeurs}
Le \textit{Dictionnaire de linguistique} de Larousse définit ainsi les embrayeurs comme une classe de mots dont le sens varie avec la situation. Ces mots, n’ayant pas de référence propre dans la langue, ont la particularité de ne recevoir un référent que lorsqu’ils sont inclus dans un message.

Si l’embrayeur a pour fonction de représenter la réalité extralinguistique, cette représentation s’opère de manière tout à fait relative, puisque l’identification de l’objet désigné ne peut être réalisée que si l’on connaît la situation d’énonciation.

\textbf{Ex :} Je vais me coucher.

Si l’on sait que c’est Paul qui a prononcé cet énoncé, on peut en déduire que « je » représente Paul. Mais si l’on ignore qui l’a prononcée,\ednote{prononcée 的前置宾语 l’ 指阳性的 énoncé，应为 prononcé。} on ne saura pas qui est dénoté par « je ».

\minititle{\textcircled{\scriptsize 1} Les embrayeurs subjectifs}
Les embrayeurs renvoient soit aux actants soit aux circonstants de l’énonciation. Les premiers, qui se réfèrent à l’énonciateur ou au destinataire, sont appelés embrayeurs subjectifs (ou personnels). Les seconds se réfèrent aux circonstances de l’énonciation : circonstances de lieu embrayeurs spatiaux),\ednote{lieu 后缺左括号，应为 circonstances de lieu (embrayeurs spatiaux)。} ou circonstances de temps (embrayeurs temporels).

\minititle{1. Embrayeurs de la première personne}
La présence de l’énonciateur peut se manifester par les embrayeurs suivants :
\begin{itemize}
\item Les marques de la première personne du singulier ou du pluriel : cela concerne la terminaison des verbes conjugués, les pronoms personnels (« je », « me », « moi », « nous », etc.), les adjectifs et pronoms possessifs (« mon », « ma », « mes », « notre », « le nôtre », etc.).
\item Les termes relationnels ou affectifs pour lesquels est sous-entendu un possessif de la première personne :
\end{itemize}
\textbf{Ex :} J’ai rencontré un ami. (sous entendu : J’ai rencontré un \textit{de mes} amis.)

\hspace*{1.5em}Grand-père a téléphoné. (sous entendu : \textit{Notre} grand-père a téléphoné.)
\begin{itemize}
\item Le mode impératif, la phrase interrogative directe et l’apostrophe témoignent de la présence de l’énonciateur.
\end{itemize}
\textbf{Ex :} Réponds-moi.

\hspace*{1.5em}Quelle heure est-il ?

\hspace*{1.5em}Antoine, il commence à pleuvoir.

\minititle{2. Embrayeurs de la deuxième personne}
La présence du destinataire se manifeste par les embrayeurs suivants :

\origpage{358}{344}
\begin{itemize}
\item Les marques de la deuxième personne du singulier ou du pluriel : cela concerne la terminaison des verbes conjugués, les pronoms personnels (« tu », « te », « toi », « vous »), les adjectifs et pronoms possessifs (« ton », « tes », « votre », « le tien », etc.).
\item les termes relationnels ou affectifs pour lesquels est sous-entendu un possessif de la deuxième personne :
\end{itemize}
\textbf{Ex :} Un voisin voudrait te voir. (sous entendu : Un \textit{de tes} voisins voudrait te voir.)

\hspace*{1.5em}Maman a téléphoné. (\textit{Notre} maman a téléphoné.)\ednote{本例明列于第二人称项下，括注却用第一人称 Notre。若母亲为双方共有，Notre 可以成立；若专举第二人称隐含领属，应为 Ta。原文情境不足，保留并注明例证与分类的不完全对应。}
\begin{itemize}
\item le mode impératif, phrase interrogative directe et l’apostrophe témoignent de la présence du destinataire.
\end{itemize}
\textbf{Ex :} Réfléchissez avant de parler.

\hspace*{1.5em}Où sont mes lunettes ?

\hspace*{1.5em}Cécile, je monte me coucher.

\minititle{3. Modalisateurs du discours}
Un modalisateur du discours est un type d’embrayeur qui permet à l’énonciateur d’exprimer vis-à-vis de son énoncé ses émotions et ses jugements. Grâce aux modalisateurs, l’énonciateur va pouvoir nuancer la vérité de son énoncé, et exprimer le degré de certitude que celui-ci lui accorde. Un modalisateur peut consister en :
\begin{itemize}
\item un choix lexical particulier. Par exemple, un mot péjoratif ou mélioratif :
\end{itemize}
\textbf{Ex :} Mon cousin a été arrêté \textit{par la police.} / Mon cousin a été arrêté \textit{par les flics.}
\begin{itemize}
\item un choix flexionnel particulier. Par exemple, l’emploi du conditionnel au lieu du présent de l’indicatif, pour exprimer le doute :
\end{itemize}
\textbf{Ex :} Le mari a assassiné sa femme / le mari \textit{aurait} assassiné sa femme\ednote{此例实际对比直陈式复合过去时 a assassiné 与条件式过去时 aurait assassiné；前文“直陈式现在时”与例句不符。}
\begin{itemize}
\item L’insertion d’un élément linguistique. On peut citer par exemple :
\begin{itemize}[label=--]
\item une interjection ou une exclamation : hélas, pauvres de nous, etc.
\item un adverbe ou une locution adverbiale : franchement, sincèrement, malheureusement, pas du tout, etc.
\item un syntagme nominal : sans doute, selon moi, d’après eux, etc.\ednote{sans doute 为副词性短语，selon moi、d’après eux 为介词性结构；不能把三个完整形式一律标为 syntagme nominal（名词短语）。}
\item une proposition incise : à dire vrai, pour parler franchement, à ce que l’on dit, si mes souvenirs sont exacts, etc.
\end{itemize}
\end{itemize}

\origpage{359}{345}
\minititle{\textcircled{\scriptsize 2} Les embrayeurs spatio-temporels}
Comme leur nom l’indique, ces embrayeurs se divisent en embrayeurs spatiaux et embrayeurs temporels. Comme nous le verrons un peu plus bas, ils sont aussi appelés « embrayeurs déictiques ».

\minititle{1. Embrayeurs spatiaux}
Nous l’avons vu : les circonstances de lieu de l’énonciation s’apprécient par rapport au lieu où le « je » parle (ou écrit). Elles se manifestent par les embrayeurs suivants :
\begin{itemize}
\item certains adverbes de lieu (« ici », « là-bas », « là-haut », « céans », « derrière », etc.)
\item certains syntagmes nominaux (« à gauche », « près de moi », etc.)\ednote{à gauche、près de moi 整体为介词性／副词性空间表达，而非名词短语；这是句法类别的误标，不影响其可具有空间指示功能。}
\item les présentatifs « voici » et « voilà »
\item les adjectifs et pronoms démonstratifs accompagnés d’un geste de monstration :
\end{itemize}
\textbf{Ex :} Parmi tous ces gâteaux, je choisis celui-ci.
\begin{itemize}
\item les verbes de déplacement : il s’agit essentiellement du verbe « venir » quand il signifie « se rapprocher du lieu de l’énonciation » et des verbes « aller » et « s’en aller » quand ils signifient s’éloigner du lieu d’énonciation.
\end{itemize}

\minititle{2. Embrayeurs temporels}
Les circonstances du temps de l’énonciation s’apprécient par rapport au moment où le « je » parle (ou écrit). Elles se manifestent par les embrayeurs suivants :
\begin{itemize}
\item certains adverbes de temps : « maintenant », « hier », « demain », « actuellement », etc.
\item certains syntagmes nominaux : « en ce moment », « le mois dernier », « depuis une semaine », etc.
\item les terminaisons des verbes principaux : présent d’énonciation, passé composé, futur, etc.
\end{itemize}
Notons que lorsqu’ils sont associés à la troisième personne, les temps du récit (présent de narration, passé simple et imparfait) ne constituent pas des embrayeurs temporels, car ils se trouvent ainsi hors discours oral. On dit qu’ils sont « hors plan embrayé ».\ednote{第三人称加叙述时态并不足以推出“脱离口头话语”或无指示定位。例如 Il dormait ici hier 即为第三人称未完成过去时，却由 ici、hier 明确联系说话情境。“非接合叙述”须按整个话语的定位方式判断。}

Dans certaines circonstances, le mot « déictique » est employé comme synonyme d’embrayeur. Au sens strict, le déictique est un embrayeur particulier, réservé à l’usage oral, souvent accompagné d’un geste de monstration.\ednote{“仅限口语”过窄：书信中的 je、ici、demain 同样可具有指示功能，且指示不必伴随手势。下文自己的“时间指示”定义也并未要求可见手势。术语的较窄用法应与一般的 deixis 概念区别。}

\subsection{Les déictiques}
La notion de deixis apparaît en 1904 chez le linguiste allemand \textbf{Karl Brugmann} (1849-1919) qui appelle déictiques les pronoms, les démonstratifs et certains adverbes de temps et de lieu qui ont la possibilité de référer directement à la situation extérieure. Le concept a ensuite été repris par \textbf{Karl\origcont{360}{346} Bühler}, puis par \textbf{Charles Bally} dans sa \textit{Linguistique générale et linguistique française} (1932).

\minititle{\textcircled{\scriptsize 1} Définition}
Deixis est un mot grec qui signifie « montrer ».\ednote{希腊语 {\BBSixGreek δεῖξις} 是名词，义为“指示／展示”；“montrer”是对相关动词 {\BBSixGreek δείκνυμι} 的释义。此处不宜把名词词性误解为动词。} Elle intervient lorsque la compréhension de certaines parties d’un énoncé nécessite une information contextuelle. On dira donc qu’un mot ou une expression est déictique si son interprétation varie en fonction du contexte, comme un pronom par exemple.

\minititle{\textcircled{\scriptsize 2} Les formes de déixis}
Très schématiquement, on peut dire que tout locuteur, en prenant la parole, établit un ensemble de trois coordonnées liées à la situation d’énonciation et manifestées par les déictiques : \textit{ego, hic, nunc} (en latin : « moi », « ici », « maintenant »).

\minititle{1. La déixis de personne}
La première consiste en une référence à la personne qui parle ou aux personnes en présence au moment de l’énonciation. Elle se matérialise avec les déictiques de personne tels que « je » et « tu ». Lorsque des conventions sociales s’appliquent, avec par exemple le vouvoiement en français, on parle de déixis sociale, terme qui dénote la façon dont les statuts sociaux se reflètent dans les mots utilisés. Plus précisément, on parle de « déixis sociale relationnelle » pour les exemples du tutoiement et du vouvoiement qui dépendent des relations sociales entre les deux interlocuteurs, et de « déixis sociale absolue » pour des exemples tels que « Sa Majesté » ou « Monsieur le président », où seul le statut de l’interlocuteur compte.

\minititle{2. La déixis spatiale}
La déixis spatiale est une référence à un élément visible du lieu de l’énonciation, éventuellement au lieu même de l’énonciation. Elle se matérialise avec les déictiques de lieu tels que « ici » et « là ».

Lorsqu’un contraste entre une proximité immédiate et un espace plus distant est nécessaire, on peut utiliser les marqueurs déictiques « -ci » et « -là », qui apparaissent en particulier dans les pronoms démonstratifs « celui-ci » et « celui-là », mais qui sont susceptibles d’apparaître dans n’importe quel groupe nominal démonstratif (« cet objet-ci », « cet objet-là »).

\minititle{3. La déixis temporelle}
La déixis temporelle est en une référence au moment de l’énonciation.\ednote{est en une référence 应为 est une référence，或 consiste en une référence；两种结构发生混合。} Elle se matérialise avec les déictiques de temps tels que « maintenant », « hier » et « tout de suite ».

Pour chacune de ces trois formes de déixis, un geste peut renforcer l’ancrage dans le contexte d’énonciation, mais c’est dans le cas de la déictique spatiale qu’il est le plus utile lorsqu’il aide à l’identification d’un référent parmi plusieurs candidats possibles.

\minititle{4. Exemple}
Prenons un exemple pour illustrer cela.

\origpage{361}{347}
\textbf{Ex :} Je reviendrai demain.

Dans cet énoncé, on trouve trois éléments déictiques :
\begin{itemize}
\item Le plus apparent est le pronom personnel « je », qui apparaît aussi dans la flexion verbale « -ai ». Pour savoir qui est désigné par « je », pour identifier cette « première personne », il faut savoir qui prononce l’énoncé. Ce renseignement est normalement fourni par la situation d’énonciation : l’auditeur entend et voit (mais pas dans l’obscurité ni au téléphone) la personne qui parle et qui lui est ainsi « montrée » par la situation d’énonciation, d’où le terme de deixis. Le déictique « je » invite donc l’auditeur à compléter le sens en le reportant à la situation.
\item Le second déictique de l’énoncé est le morphème de futur « -r- ». Mais c’est un futur \textit{relatif} à la situation d’énonciation, qui est l’instant où l’énonciateur est en train de parler.
\item Le troisième déictique est « demain ». Nous avons vu dans le chapitre 9 que « demain » est un morphème dépendant. En effet, dans une situation de communication, « demain » sous-entend toujours « demain \textit{par rapport au moment où je parle.} » C’est un déictique temporel exprimant un futur \textit{relatif} au moment de la situation d’énonciation.\ednote{这里“dépendant”指指称解释依赖语境，不等于形态学上必须附着于另一形式的黏着语素。demain 可以独立成答语，应区分语境依赖与形态依附。}
\end{itemize}

\subsection{Endophore et exophore}
Nous l’avons vu en début de ce chapitre : linguistiquement, un objet peut avoir deux lieux d’existence : dans le discours ou hors du discours. Ce postulat a permis au linguiste \textbf{Michael Halliday} et à son épouse \textbf{Ruqaiya Hasan} en 1976 d’établir une différence entre endophore (appelée « référence textuelle »), lorsque le référent se trouve dans l’espace textuel, et exophore (« référence situationnelle ») lorsque le référent de l’expression se trouve localisé dans l’espace non discursif.

\minititle{\textcircled{\scriptsize 1} L’endophore}
L’endophore comprend l’anaphore et la cataphore. En linguistique, l’étude de ce qu’on appelle « les marques d’ancrage référentiel de l’énoncé » n’a cessé d’inspirer à divers auteurs une riche terminologie : en 1989, le linguiste français \textbf{Michel Maillard} a regroupé anaphore et cataphore sous le terme de diaphore.

\begin{lecture}{}
Après avoir utilisé quelque temps le mot anaphore comme le terme non-marqué de l’opposition anaphore / cataphore, nous avons ressenti le besoin d’utiliser un hyperonyme pour « couvrir » ces deux processus de sens contraire.
\end{lecture}

Dans les lignes qui vont suivre, on va rester fidèle à la terminologie de Halliday.

\minititle{1. L’anaphore}
L’anaphore est un procédé consistant à utiliser un élément du discours (le plus souvent un pronom) pour renvoyer à un constituant qui le précède (son antécédent) afin d’en permettre l’identification et l’interprétation.

\origpage{362}{348}
\textbf{Ex :} J’ai croisé mon voisin hier : il avait mauvaise mine.

Dans l’exemple ci-dessus, « il » se rapporte à « mon voisin » et ne peut s’interpréter qu’en rapport à ce dernier : on parle donc de relation anaphorique ou de « déictique anaphorique ».

\minititle{1) Définition}
\textbf{Oswald Ducrot} et le sémiologue \textbf{Tzvetan Todorov} (1939-2017) définissent l’anaphore en termes d’interprétation et considèrent qu’un segment de discours est anaphorique s’il faut se reporter à une autre partie de ce même discours pour lui donner une interprétation. Dans cette optique, l’anaphore est donc un phénomène de dépendance interprétative de deux unités dont la première, à laquelle se reporte la seconde (anaphorique), est appelée « interprétant ».

\minititle{2) Fonctions}
Du point de vue fonctionnel, l’anaphore permet d’éviter une répétition.

\textbf{Ex :} Jean n’avait pas de stylo : je lui ai prêté \textit{le mien}

Le pronom possessif « le mien » est une anaphore dont l’antécédent est le « stylo », et qui évite la répétition de ce dernier.

Une autre fonction de l’anaphore est de dissiper une amphibologie (une équivoque).

\textbf{Ex :} Mon voisin a adopté un gros chien. Il est vraiment très laid.

Ici, le pronom personnel « il » est une anaphore, mais quel est son antécédent ? Qui est vraiment très laid ? Est-ce le voisin ou le gros chien ? Pour dissiper cette amphibologie, l’utilisation d’un syntagme nominal anaphorique (« cet animal ou « ce dernier »)\ednote{cet animal 后漏闭引号，应为 (« cet animal » ou « ce dernier »)。此外 ce dernier 倾向指较近先行项，并非在任何语境中都能机械地消除歧义。} permettra de lever le doute sur l’identité de son antécédent.

\minititle{3) Catégories}
Trois catégories principales peuvent jouer le rôle d’anaphore : le nom, le pronom, et l’adverbe.

\textbf{Ex :} J’ai adopté un petit chien : \textit{cet animal} est très affectueux.

\hspace*{1.5em}Cet étudiant n’avait pas de livre : je lui ai prêté \textit{le mien.}

\hspace*{1.5em}J’ai invité ma mère à venir à la maison. \textit{Ici,} elle sera au calme.

\minititle{2. La cataphore}
Figure inverse de l’anaphore, la cataphore est un mot (ou un syntagme) qui, dans un énoncé, précède et renvoie sémantiquement à un segment à venir, appelé conséquent :

\textbf{Ex :} Quand tu \textit{le} verras, tu diras au professeur que je déteste la linguistique.

Le pronom personnel « le » est une cataphore. Son conséquent est le syntagme nominal « le\origcont{363}{349} professeur ».

La première mention de la cataphore apparaît là encore chez \textbf{Karl Bühler}, qui établit une distinction entre « monstration anticipante » (\textit{Vorverweis, en allemand}) et « monstration rétrospective », (\textit{Rückverweis}) et propose d’utiliser deux autres termes d’origine grecque : \textit{kataphora} et \textit{anaphora}. Plus tard, dans sa \textit{Linguistique générale et linguistique française} (1950), \textbf{Bally} utilise le terme d’« anticipation » pour remplacer le terme de cataphore, et définit l’anticipation comme un fait qui se produit chaque fois qu’« un signe nécessaire à la compréhension d’un autre précède celui-ci au lieu de le suivre ». Plus proche de nous chronologiquement mais plus éloigné géographiquement, le terme de cataphore a été remplacé par celui de « pronominalisation régressive » (\textit{backwards pronominalization}) par les linguistes générativistes américains, tandis que l’anaphore apparaît sous le nom de pronominalisation progressive (\textit{forwards pronominalization}). Notez que dans une énonciation, anaphore et cataphore peuvent coexister, comme dans l’énoncé suivant :

\textbf{Ex :} Quand tu \textit{le} verras, tu diras au professeur que le livre de linguistique qu’\textit{il} a écrit me donne des cauchemars.

Dans l’exemple ci-dessus, le pronom personnel « le » est une cataphore du syntagme nominal « le professeur », et le pronom personnel « il » est une anaphore de ce même syntagme nominal. Le représenté, c’est-à-dire, le professeur, est à la fois le conséquent (du pronom personnel « le ») et antécédent (du pronom personnel « il »).

\minititle{\textcircled{\scriptsize 2} L’exophore}
Quand le référent de l’expression se trouve localisé hors de l’espace discursif, \textbf{Halliday} et \textbf{Hasan} parlent d’exophore (référence situationnelle). Les linguistes \textbf{Thomas Fraser} et \textbf{André Joly} en distinguent deux types : l’exophore mémorielle (« deixis \textit{in absentia} ») et l’exophore a-mémorielle (« deixis \textit{in praesentia} »).

\minititle{1. Exophore mémorielle}
La déixis mémorielle se caractérise par le recours à des connaissances, des références ou à des événements extralinguistiques partagés entre énonciateur destinataire,\ednote{énonciateur destinataire 之间缺连接词 et；应为 entre énonciateur et destinataire。} et reconstruits dans l’esprit de l’interlocuteur. Les exophores mémorielles font appel à la mémoire de l’interlocuteur (d’où le nom de « mémorielles ») puisque le référent n’est pas présent dans le contexte (d’où le nom de déixis « \textit{in absentia} »).

\textbf{Ex :} Tu te souviens de cet orage le jour de notre mariage ?

Ici, l’énonciateur rappelle au destinataire leur mariage, sans doute relativement éloigné dans le temps. Le démonstratif « cet » ne doit pas induire en erreur : il n’est là que pour témoigner que par la pensée, il est toujours présent dans la situation d’énonciation et fait partie des souvenirs communs aux deux actants en présence, à savoir le mari et la femme.

La deixis in absentia peut prendre diverses formes, parmi lesquelles on peut citer :

\origpage{364}{350}
\begin{itemize}
\item La référence ostensive gestuelle indirecte : le locuteur montre un objet présent dans la situation d’énonciation pour indiquer à l’interlocuteur le véritable référent. Par exemple : pointant du doigt une copie d’examen sur laquelle est écrite la note de zéro sur cent, un professeur s’écrie : « \textit{Cet} étudiant n’a pas révisé l’examen ! » Le geste de monstration désigne la copie pour suggérer que le véritable référent est l’étudiant. Cela consiste à montrer quelque chose pour désigner autre chose d’absent dans la situation d’énonciation. Autre exemple : « Cet automobiliste a dû être pressé, en désignant une voiture vide garée au milieu de la route. »
\item La référence démonstrative non gestuelle indirecte. Dans une salle d’attente, deux patientes attendent le même médecin. Sans mention antérieure de ce dernier, l’une des deux dit à l’autre : « il paraît que \textit{ce} docteur est très doux avec ses patients ». Autre exemple : « ce train a toujours du retard », prononcé sur le quai de la gare sans mention antérieure par un locuteur qui attend le même train que son interlocuteur.
\item La référence générique, où la présence d’occurrences spécifiques permet de référer à la classe générique à laquelle elles appartiennent. Elle peut être soit démonstrative (« Paul a acheté une Toyota : ces voitures sont robustes. »), soit pronominale (« En Italie, \textit{ils} roulent comme des fous. »). \textit{Autre exemple : « Attention ! Ils sont dangereux. » (prononcé en présence d’un} référent spécifique, un chien, par exemple, pour avertir du danger que représentent les chiens en général).
\item L’emploi du pronom personnel « il » sans antécédent pour viser un référent non présent. Par exemple, « il va vous recevoir tout de suite ! », prononcé par la secrétaire aux deux patientes qui attendent le docteur dans la salle d’attente.
\end{itemize}

\minititle{2. Exophore a-mémorielle}
La deixis \textit{in praesentia} (exophores a-mémorielles) ne fait pas appel à la mémoire puisque le référent est présent dans le contexte : je peux montrer du doigt (ou du regard) le référent. La deixis \textit{in praesentia} fait emploi de la situation immédiate. Voici un dialogue téléphonique impliquant une deixis \textit{in praesentia} :

\textbf{Client :} Bonsoir, je souhaite réserver une table pour deux, ce soir, à 19 heures.

\textbf{Serveur :} (montrant du doigt une table) : Bien sûr. Celle-ci, ou bien celle-là ?

Nous avons ici un cas d’exophore a-mémorielle, mais celle-ci ne peut pas être perçue par le client puisqu’il ne peut pas percevoir le signe de monstration.

\minititle{\textcircled{\scriptsize 3} Hétérogénéité des marqueurs référentiels}
En principe, l’anaphore et la déixis sont opposées. En effet, une expression anaphorique qui marque avant tout la continuité avec un référent déjà placé dans le focus, alors qu’une expression déictique a précisément pour rôle d’attirer l’attention de l’interlocuteur sur un nouvel objet de référence.\ednote{原句 une expression anaphorique qui marque… alors que… 缺少主句谓语；按与后半句平行的结构，qui 宜删去，作 une expression anaphorique marque…。} En fait, les choses ne sont pas si tranchées que cela : aucune expression n’est uniquement anaphorique ou déictique. Le contexte sert de situation : certains déictiques peuvent\origcont{365}{351} indifféremment servir à montrer ce qu’on a sous les yeux dans la réalité (deixis in praesentia), ou à renvoyer à des mots du contexte (emploi anaphorique de la déixis). Par conséquent, l’article défini, l’adjectif et le pronom démonstratif, l’adjectif possessif, et le pronom personnel de la troisième personne connaissent deux types d’emplois illustrés dans les énoncés ci-dessous :

\textbf{Ex :} Fais attention à \textit{la} marche ! (emploi de situation immédiate)

\hspace*{1.5em}Paul glissé sur une marche.\ednote{Paul glissé 缺助动词，应为 Paul a glissé。} \textit{La} marche était mouillée après la pluie. (emploi anaphorique)

\textbf{Ex :} \textit{Cette} voiture a le tuyau d’échappement percé. (avec geste d’ostension sur la voiture)

\hspace*{1.5em}Paul a acheté une voiture. Il sait que \textit{cette} voiture lui simplifiera la vie. (emploi anaphorique)

\textbf{Ex :} Je veux \textit{celle-ci.} (emploi de situation immédiate en pointant une pâtisserie)

\hspace*{1.5em}J’ai interrogé ma fille sur le vase cassé. \textit{Celle-ci} ne sait rien. (emploi anaphorique)

\textbf{Ex :} Ne te brûle pas main !\ednote{pas main 应为 pas la main，缺少定冠词。} \textit{Il} sort du four ! (emploi de \textit{il} sans antécédent, montrant un gâteau)

\hspace*{1.5em}Cet étudiant a assassiné son professeur. \textit{Il} déteste la linguistique. (emploi anaphorique)

C’est sur cet exemple, qui, j’espère, ne vous donnera pas l’idée de faire la même chose, pas,\ednote{chose 后的 pas 游离于句法结构，疑为残留误植；可删去，保留其余句子。} que s’achève la partie sur la linguistique de l’énonciation. Nous allons, dans la dernière partie de ce chapitre, nous pencher sur les prolongements récents de la linguistique de l’énonciation, essentiellement en Angleterre et aux États-Unis : il s’agit de la linguistique pragmatique. Nous verrons que parler, c’est agir.

\section{LA PRAGMATIQUE : DIRE, C’EST FAIRE}
Héritage de la rhétorique gréco-latine dont elle a théorisé les conceptions, la pragmatique est la branche de la linguistique qui étudie l’emploi de la langue en situation. Son postulat de base est que le sens n’est pas seulement inscrit dans les énoncés eux-mêmes : il se construit aussi dans l’activité des sujets engagés dans une communication, elle-même inscrite dans un contexte d’emploi particulier.

C’est le sémiologue américain \textbf{Charles Morris}, présenté dans le chapitre 6, qui a forgé le terme de pragmatique en 1938 et a défini cette nouvelle discipline comme l’étude des relations qui existent entre les signes et les usagers de ces signes.\origfoot{2}{« \textit{the study of the relation of signs to interpreters} », cité dans Morris (Charles W.) \textit{Foundations of the Theory of Signs}, Chicago, The University of Chicago Press., 1938, 6p.}

Dans un ouvrage postérieur, il précise :
\begin{lecture}{}
Pragmatics is that portion of semiotic which deals with the origin, uses, and effects of signs within the behavior in which they occur.
\end{lecture}

\origpage{366}{352}
Dans les années 60 et 70, cette discipline récente a connu un essor sous l’effet d’une part des travaux du philosophe anglais \textbf{John Langshaw Austin} (1911-1960) à Oxford, et d’autre part des travaux de Benveniste et Jakobson sur les embrayeurs. Si des divergences apparaissent entre linguistes du fait de la grande diversité de leurs travaux, ils partagent néanmoins une même conception du langage et s’opposent à toute dissociation d’une phrase et de son sens. Pour les pragmaticiens, l’énonciation n’est pas quelque chose d’accessoire qui s’ajoute à une phrase déjà porteuse de sens en elle-même : l’énonciation conditionne radicalement le sens de la phrase. Dans cette conception, l’idée très longtemps dominante en philosophie selon laquelle le langage n’a pour fonction que de représenter le monde et de transmettre des informations est mise à mal : le langage est un acte. Pour le linguiste américain \textbf{John Searle} (1932-), « élève et continuateur de John Austin, « Parler une langue, c’est réaliser des actes de langage ».\ednote{élève 前的开引号没有对应闭引号；宜删去该开引号，将引号仅用于后面的直接引语。生卒年括注按底本保留，不作为当前人物状态资料。} Son célèbre ouvrage éponyme, \textit{The acts of language} (1970),\ednote{所列 The acts of language 不是英文原著书名；英文原著为 Speech Acts: An Essay in the Philosophy of Language（1969）。依据英文书目记录；法译本版年另待核。\ecite{58}} est un des ouvrages fondateurs cette nouvelle discipline.\ednote{fondateurs cette nouvelle discipline 中缺 de，应为 fondateurs de cette nouvelle discipline。}

\subsection{Les énoncés constatifs et performatifs}
Professeur de philosophie morale à Oxford, John Langshaw Austin a donné une série de conférences à l’université de Harvard en 1955 qui ont été publiées en 1962 sous le titre \textit{How to do Things with Words}. Le livre sera traduit en français en 1970 sous le titre « Quand dire, c’est faire ». Le point de départ d’Austin est que les philosophes ont longtemps supposé qu’une affirmation ne pouvait que \textit{décrire} un état de fait :

\textbf{Ex :} Hier, il a plu.

\hspace*{1.5em}Je regrette d’avoir choisi le français comme spécialité.

Ces énoncés, appelés constatifs, sont soit vrais soit faux. L’apport d’Austin consiste à avoir montré qu’il existe des énoncés qui ont la particularité d’être eux-mêmes l’acte qu’ils désignent, et qui donc n’entrent pas dans la catégorie des énoncés constatatifs :\ednote{constatatifs 应为 constatifs，多植 ta。} ce sont les énoncés performatifs, qui ne décrivent pas la réalité, mais agissent sur elle.

\textbf{Ex :} Que la lumière soit !

La différence entre les énoncés performatifs et les énoncés constatifs réside dans ce qu’on appelle la « direction d’ajustement ». Les énoncés constatifs ont pour but de décrire le réel, donc de s’ajuster à lui. Le réel reste, après l’émission de l’énoncé, ce qu’il était auparavant. À l’opposé, après un énoncé performatif, le réel n’est plus tout à fait ce qu’il était auparavant. Dans ce cas, c’est le réel qui s’ajuste à l’énoncé.\ednote{“适配方向”不能与“施为／记述”一一对应：Searle 的分类中，断言为词适配世界，指令与承诺为世界适配词，表达类无此适配方向，宣告类为双向。施为行为也不保证期望的外部结果已经发生。\ecite{59}}

\minititle{\textcircled{\scriptsize 1} La notion de performativité}
Cherchant à décrire et à nommer ce type particulier d’énonciations (\textit{utterances}) qui, comme les promesses, les souhaits ou les jugements, ne se contentent pas de dire mais font par la même occasion, Austin définit ainsi la performativité :

\origpage{367}{353}
\begin{lecture}{}
Le terme performatif dérive, bien sûr, du verbe anglais perform : il indique que produire l’énonciation revient à exécuter une action.» Prononcer des mots, en effet, est d’ordinaire un événement ou même l’événement principal, dans le fait d’accomplir (perform) l’acte.\origfoot{3}{Texte original : « The term […] « performative » is derived, of course, from “perform” […] : it indicates that the issuing of the utterance is the performing of an action. The uttering of the words is, indeed, usually a, or even the, leading incident in the performance of the act ».}
\end{lecture}

On voit que la performativité est le fait, pour un signe linguistique (énoncé, phrase, verbe, etc.) de réaliser lui-même ce qu’il énonce, comme dans l’énoncé ci-dessous, prononcé par le maire lors d’un mariage :

\textbf{Ex :} Je vous déclare mari et femme.

Notez toutefois que l’énonciation n’est performative que si les différents protagonistes respectent certaines conditions de succès que Austin appelle \textit{felicity conditions} (« conditions de félicité ») : le locuteur doit être maire, les destinataires célibataires, etc. Ainsi, le même énoncé, prononcé par un simple employé de mairie ne fera pas changer les deux fiancés de statut. Il faut donc prêter une certaine attention au contexte dans lequel se produit l’énonciation. C’est ce qu’illustre Austin dans cet extrait de \textit{How to do things with words} (1962) :

\begin{lecture}{}
Supposons, par exemple, que j’aperçoive un bateau dans une cale de construction, que je m’en approche et brise la bouteille suspendue à la coque, que je proclame « Je baptise ce bateau Joseph Staline », et que, pour être bien sûr de mon affaire, d’un coup de pied je fais sauter les cales. L’ennui, c’est que je n’étais pas la personne désignée pour procéder au baptême (peu importe que Joseph Staline ait été ou non le nom prévu - ce ne serait qu’une complication de plus. Nous admettons sans peine :

1) que le bateau n’a pas, de ce fait, reçu de nom ;

2) qu’il s’agit d’un incident extrêmement regrettable.

On pourrait dire que j’ai « rempli certaines formalités » de la procédure destinée à baptiser le bateau, mais que mon « action » fut « nulle et non avenue » ou « sans effet », parce que je n’étais pas la personne adéquate, que je n’avais pas les « pouvoirs » pour l’accomplir. Mais on pourrait dire aussi […] que lorsqu’il n’y a ni prétention ni même ombre d’un droit aux pouvoirs, alors il n’existe aucune procédure conventionnelles reconnue : « c’est une imitation bouffonne, comme un mariage avec un singe. »

[…] Il est toujours nécessaire que les circonstances dans lesquelles les mots sont prononcés soient d’une certaine façon (ou de plusieurs façons) appropriées, et qu’il est d’habitude nécessaire que celui-là même qui parle, ou d’autres personnes, exécutent aussi certaines autres actions - actions « physiques » ou « mentales », ou même actes consistant à prononcer ultérieurement d’autres paroles. C’est ainsi que pour baptiser un bateau, il est essentiel que je sois la personne désignée pour le faire ; que pour me marier (chrétiennement), il est essentiel\origcont{368}{354} que je ne sois pas déjà marié avec une femme vivante, saine d’esprit et non divorcée, etc. Pour qu’un pari ait été engagé, il est nécessaire en général que la proposition du pari ait été acceptée par un partenaire (lequel a dû faire quelque chose, dire « D’accord! », par exemple). Et l’on peut difficilement parler d’un don si je dis « Je te le donne », mais ne tend point l’objet en question.\ednote{长引文有三处可明确定位的排印／语法问题：plus 后缺闭括号；procédure conventionnelles 应为 procédure conventionnelle；je… ne tend 应为 je… ne tends。均按底本照录，不据此反推为 Austin 原著的错误。}
\end{lecture}

\minititle{\textcircled{\scriptsize 2} Les expressions de la performativité}
La performativité d’un énoncé peut bien sûr être exprimée au moyen d’un verbe performatif, mais il existe d’autres manières.

\minititle{1. Les énoncés à performativité lexicalement dénommée}
Les énoncés à « performativité lexicalement dénommée » représentent le cas le plus fréquent : ce sont ceux qui comprennent un verbe performatif :

\textbf{Ex :} Je vous \textit{promets} de ne plus arriver en retard en cours.

\hspace*{1.5em}Je vous \textit{déclare} coupable du meurtre de votre professeur de linguistique.

\hspace*{1.5em}\textit{Je vous nomme délégué de classe.}

\hspace*{1.5em}Je vous \textit{pardonne} d’avoir encore oublié la fête de professeurs.

\minititle{2. Les énoncés à performativité indiquée autrement}
La « performativité indiquée autrement » est une autre façon de dire « sans verbe performatif ». Dans ce cas, la performativité peut consister en divers procédés grammaticaux, comme :
\begin{itemize}
\item Le mode impératif : Ouvrez votre livre à la page 123.
\item Le mode interrogatif : Tu ne finis pas ta pizza ?
\item Les interjections et expressions toutes faites : Chut ! Halte ! Bis ! Ouste ! Du balai ! Stop ! Salut ! Bonjour ! Merci ! Objection ! Pardon ? Hein ? Comment ?
\end{itemize}

\minititle{\textcircled{\scriptsize 3} Limites de la distinction entre performatif et constatif}
Dans un énoncé, l’adverbe « dehors » peut être descriptif s’il est une réponse à une question. Mais il peut aussi être performatif s’il s’agit d’un ordre bref accompagné d’une intonation injonctive. En plus des performatifs explicites vus plus haut, il existe aussi des performatifs implicites dans lesquels la performativité, bien que non exprimée, est bien réelle. Comparons les deux exemples ci-dessous :

\textbf{Ex :} \textit{Je promets} d’étudier plus sérieusement le semestre prochain.

\hspace*{1.5em}J’étudierai plus sérieusement le semestre prochain.

Le deuxième énoncé est une forme de promesse ou d’engagement, même si le verbe promissif\origcont{369}{355} « promettre » n’est pas explicitement employé. Prenons un autre exemple :

\textbf{Ex :} Le directeur est dans son bureau.

Prononcé par une secrétaire s’adressant à un employé, cet énoncé en apparence constatif modifie en fait la réalité. Il a enrichi les connaissances du destinataire et va probablement susciter une réaction chez lui, à savoir se rendre dans le bureau du directeur. Cet énoncé, tout en restant constatif, a bien une dimension performative puisque l’énonciation fait d’une pierre deux coups : décrire, et informer.

Austin constate qu’une énonciation unique a des effets multiples. Il note également que les énoncés constatifs sont, comme les performatifs, soumis à des conditions de félicité et donc eux aussi sujets aux échecs et aux ratages. Par exemple, la promesse implicite d’étudier plus sérieusement le semestre prochain ne sera peut-être pas tenue. De même, l’employé ne comprendra peut-être pas qu’il est convoqué chez le directeur. Austin en conclut que l’opposition initiale entre les conditions de félicité des énoncés performatifs et les conditions de vérité énoncés constatifs\ednote{vérité 与 énoncés 之间缺 des，应为 les conditions de vérité des énoncés constatifs。} n’est plus tranchée, puisqu’en effet elles peuvent se combiner dans le même énoncé, comme dans l’exemple suivant, que Galilée aurait pu prononcer :

\textbf{Ex :} J’affirme que la Terre est ronde.

La distinction initiale entre performatif et constatif semble se brouiller : ce qui vaut pour les performatifs vaut pour les constatifs, notamment en ce qui concerne les conditions de félicité, car ils fonctionnent en réalité de la même façon, même s’ils ne servent pas à faire la même chose. Autrement dit, les énoncés descriptifs sont également des actes, comme les énoncés performatifs : ce sont tous des « actes de langage ».

\subsection{Les actes de langage}
Généralisant la performativité à tous les énoncés, Austin va dresser une nouvelle typologie des énoncés.

\minititle{\textcircled{\scriptsize 1} La classification des \textit{speech acts}}
La classification des actes de langage comprend : les actes locutoires, les actes illocutoires, et les actes perlocutoires.

\minititle{1. Les actes locutoires}
Dans la théorie linguistique de John Austin, les actes locutoires correspondent à l’exercice de la faculté de langage par un locuteur individuel. Produire un énoncé, c’est réaliser des actes locutoires, c’est-à-dire lier des phonèmes, des morphèmes, des syntagmes, concevoir des phrases, leur attribuer du sens, les prononcer (ou les écrire).

\minititle{2. Les actes illocutoires}
La seconde catégorie est celle des actes « illocutoires », c’est-à-dire des actes contenus dans le\origcont{370}{356} langage. Décrire n’est qu’une des activités que permet le langage. Les énoncés illocutoires ont une valeur d’action au-delà de leur sens immédiat. L’énoncé « Est-ce que vous avez l’heure ? » n’a pas pour fonction de savoir si l’interlocuteur possède une montre, mais exprime plutôt la volonté que ce dernier nous donne l’heure : c’est la fonction illocutoire de cet énoncé, et c’est dans la volonté de réaliser certains actes que réside la force illocutoire des énoncés. Cette volonté, cet élan, c’est, rappelez-vous, la fonction conative du langage.\ednote{施事力不能等同于“意动功能”：断言、感谢、宣告等也具有施事力，却不能全部还原成促使听者行动的功能。间接请求只是施事行为的一例。\ecite{59}}

\minititle{3. Les actes perlocutoires}
La troisième catégorie, celle des actes perlocutoires, correspond à la phase de décodage de l’énoncé. Il s’agit ici de comprendre le sens immédiat (de l’acte locutoire) d’un énoncé, mais aussi la force illocutoire (l’intention) qu’il contient, et de réagir en conséquence.\ednote{取效行为并非“解码阶段”：Austin 区分听者对意义及施事力的理解（uptake）与说服、惊吓、促使行动等进一步效果。理解请求而不照办仍可完成理解，但未产生说话者预期的取效结果。\ecite{60}} Par exemple, si au cours d’un repas je demande à mon voisin de table s’il a du sel (acte illocutoire), il est probable que ce dernier me tendra la salière : c’est l’effet perlocutoire. Mais il se peut aussi qu’il ne comprenne pas le sens véritable de ma fausse interrogation, et qu’il se contente de me répondre : « oui ». Dans ce cas là, l’acte perlocutoire a échoué.

\minititle{4. Différence entre actes illocutoires et perlocutoires}
Il est parfois délicat de distinguer les actes illocutoires des actes perlocutoires. On peut s’accorder sur le fait que les actes illocutoires sont des actes de langage qui ne peuvent échouer, justement parce qu’ils sont inséparables du langage : selon une formule familière aux pragmaticiens, « dire c’est faire ».\ednote{“施事行为不能失败”不成立，且与本章原书 pp.353、355 关于合适条件及失误的说明冲突。说出句子不保证成功实施命名、宣告等行为；例如无权限者实施命名可以无效。须区分行为是否成立与成立后承诺是否履行。\ecite{60}} Autrement dit, il suffit d’avoir dit pour avoir fait. Ainsi, une promesse est constituée dès que l’on a prononcé l’énoncé « je promets », et il faut alors la distinguer de son exécution ou de son éventuel échec. C’est ce que le philosophe français Jacques Derrida appelle la plénitude des énoncés performatifs.

\minititle{\textcircled{\scriptsize 2} La performativité selon Derrida}
Le philosophe français \textbf{Jacques Derrida} (1930-2004) est connu, entre autres choses bien sûr, pour avoir polémiqué avec John Searle, le disciple et continuateur d’Austin. Dans \textit{Marges de la philosophie} (1972), le philosophe a mis en cause la pensée de Searle et Austin en critiquant les limites imposées au concept de performativité par le contexte de réalisation. Derrida envisage au contraire la performativité comme un « pouvoir-dire et vouloir-dire absolument plein et maître de lui-même ».\ednote{引文归属与论证方向待核：本段把“绝对充盈并掌握自身”作为 Derrida 的正面主张，但紧接着的引文讨论的是施为话语产生或改变情境；仅凭这两段不能推出行为不受语境制约。需以 « Signature événement contexte » 原篇核对被引短语的主语及上下文，此处为待核判断。\ecite{61}} Il a bien sûr constaté que l’action en jeu dans les locutions performatives aboutit fréquemment à un échec, à sa non-réalisation : conformément à ce qu’avait observé Austin, très souvent l’énonciation performative échoue et échappe à son objectif. En réaction à cet échec de la performativité énonciative, Derrida insiste sur la plénitude de l’énoncé performatif :

\begin{lecture}{}
L’énoncé performatif n’a pas son référent […] hors de lui, ou en tout cas avant lui et en face de lui. Il produit ou transforme une situation, il opère ; […] cela constitue sa structure interne, sa fonction ou sa destination ».
\end{lecture}

\origpage{371}{357}
Ce renversement qu’installe Derrida au sein du concept est d’une importance cruciale puisqu’il établit que l’acte performatif se suffit à lui-même, et que ce ne sont pas les événements extérieurs ou le contexte qui légitiment sa valeur. L’action contenue dans l’énoncé performatif se suffit à elle-même, qu’elle réussisse ou qu’elle échoue. La possibilité d’échec n’est pas exclue de la réflexion, mais celle-ci ne fragilise plus le concept puisque le résultat des actions performatives n’est plus une condition de leur reconnaissance comme performatives. La réflexion de Derrida marque un tournant certain dans l’évolution du concept de performativité dans la mesure où il affirme que l’action contenue dans l’énoncé performatif peut être effective, mais peut aussi ne pas l’être.\ednote{此段的推论仍待原篇核对：“最终结果不决定是否构成施为”并不逻辑蕴含“无须语境、自给自足”。后一命题比前一命题更强，也不能直接由前页“产生或改变情境”的引文推出。保留原文，不凭摘要重构 Derrida 的完整立场。\ecite{61}}

\minititle{\textcircled{\scriptsize 3} La classification des actes illocutoires}
En 1962, doutant de la pertinence de la distinction entre énoncés constatifs et performatifs, Austin a admis que toute énonciation correspond à l’accomplissement d’un acte illocutoire et a classé les actes illocutoires en cinq grandes catégories :\ednote{1962 是 Austin 讲稿的身后出版年，不是他在世时作出这一转变的年份。该书依据 1955 年 William James 讲座整理；以下分类见第十二讲。\ecite{60}}
\begin{itemize}
\item Les verdictifs permettent d’acquitter, de condamner, de décréter, de considérer comme, de décrire, d’analyser, d’estimer, de classer, d’évaluer, de caractériser, etc.
\item Les exercitifs permettent de commander, ordonner, pardonner, léguer, plaider pour, supplier, recommander, implorer, conseiller, nommer déclarer une séance ouverte, proclamer, etc.\ednote{列举中 « nommer déclarer » 之间缺少分隔标点，应作 « nommer, déclarer »。}
\item les commissifs engagent le locuteur à une suite d’actions déterminée : promettre, faire le vœu de, s’engager, garantir, jurer, passer une convention, embrasser une cause, etc.
\item Les comportatifs : il s’agit de s’excuser, remercier, féliciter, souhaiter la bienvenue, critiquer, bénir, maudire, porter un toast, boire à la santé, protester, défier, déplorer, mettre au défi, etc.
\item Les expositifs : ils sont utilisés pour exposer des conceptions, conduire une argumentation, clarifier l’emploi des mots (affirmer, nier, répondre, objecter, concéder, exemplifier, paraphraser, rapporter des propos, affirmer, nier, postuler, remarquer, etc.)
\end{itemize}

\minititle{\textcircled{\scriptsize 4} La classification des verbes performatifs}
Entré en 1952 à l’Université d’Oxford, le linguiste américain John Searle a suivi les cours de John Langshaw Austin avant d’être nommé en 1959 professeur à l’université de Berkeley. Dans son premier ouvrage, \textit{Speech Acts} (1969), s’appuyant sur la théorie des actes de langage de son ancien professeur, il a proposé sa propre classification des verbes performatifs :\ednote{下列五类对应 Searle 后来系统提出的施事行为分类，参见 1975 年 « A Taxonomy of Illocutionary Acts »；不宜全部归于 1969 年首著。该文还明确区分“施事行为的分类”和“施事动词的分类”。\ecite{59}}
\begin{itemize}
\item Un verbe assertif engage le locuteur sur la véracité d’une proposition. Les mots s’ajustent au monde : informer, affirmer, garantir, etc.
\item Un verbe directif correspond à la tentative de la part du locuteur d’obtenir quelque chose de son destinataire. Le monde s’ajuste aux mots : demander, requérir, etc.
\item Un verbe promissif engage le locuteur sur le déroulement d’une action future. Le monde s’ajuste aux mots : promettre, garantir, assurer, s’engager, etc.
\end{itemize}
\origpage{372}{358}
\begin{itemize}
\item Un verbe expressif exprime l’état psychologique du locuteur sans direction d’ajustement : remercier, féliciter, déplorer, regretter, etc.
\item Un verbe déclaratif modifie un état institutionnel. La direction d’ajustement est double : déclarer la guerre, nommer, baptiser, etc.
\end{itemize}
Les actes illocutoires et perlocutoires évoqués plus haut nous ont permis d’introduire les notions d’échec et de ratage. Nous allons retrouver ces notions dans le point suivant, consacré à la linguistique interactionnelle.

\subsection{La pragmatique conversationnelle}
Dans une situation d’interaction verbale, il existe plusieurs situations d’échec : les silences ; les \textit{interruptions, une prise de parole trop longue, les amphibologies, etc. Pour réussir} leurs actes de langage, locuteur et interlocuteur doivent respecter des règles constitutives (apprendre à jouer) et normatives (apprendre à bien jouer). C’est ce à quoi s’intéresse la linguistique interactionniste, discipline apparue au début du 20\textsuperscript{e} siècle chez les pragmatistes américains (Charles Sanders Peirce, John Dewey et William James), mais aussi chez des intellectuels soviétiques comme Mikhaïl Bakhtine ou encore le linguiste Valentin Voloshinov (1895-1936). Tous reconnaissent, bien que de différentes manières, l’ancrage du langage dans la pratique et plus précisément dans les pratiques interactionnelles. Pour ne prendre qu’un exemple, le linguiste russe Valentin Voloshinov (1895-1936) affirmait dès 1929 :
\begin{lecture}{}
There are as many meanings of a word as there are contexts of its use. […] Word is a two-sided act. It is determined equally by whose word it is and for whom it is meant. As word, it is precisely the product of the reciprocal relationship between speaker and listener, addresser and addressee […] A word is territory shared by both addresser and addressee, by the speaker and his interlocutor.
\end{lecture}
Les années 60 ont vu différentes disciplines (sociologie, psychologie, linguistique) formuler des modèles d’analyse qui ont joué un rôle fondamental dans le développement de l’analyse de l’interaction en linguistique. Cette dernière, qui appartient à pragmatique, s’est enrichie de la théorie des actes de langage d’Austin et Searle, de la théorie de l’argumentation d’Oswald Ducrot, et de l’interactionnisme sociologique d’Erving Goffman (1922-1982).\ednote{« appartient à pragmatique » 缺冠词，应为 « appartient à la pragmatique »。}

\minititle{\textcircled{\scriptsize 1} Oswald Ducrot : les lois du discours}
Pour que les actes de langages soient efficaces, Ducrot estime que les deux parties prenant part à l’interaction verbale doivent respecter des conventions analogues aux règles d’un jeu : participer au jeu, c’est en accepter les règles. De même, utiliser le langage, c’est se soumettre à ses « lois ». Ducrot a donc formulé quatre lois du discours :\ednote{术语通常写作 « actes de langage »，此处 « langages » 的复数不合本章所用术语。}
\begin{itemize}
\item Selon la loi de sincérité, on est tenu de ne dire que ce qu’on croit vrai. Sans cette convention, aucune espèce de communication ne serait possible.
\end{itemize}
\origpage{373}{359}
\begin{itemize}
\item Selon la loi d’intérêt, on ne doit de parler que de ce qui est susceptible d’intéresser notre interlocuteur.\ednote{« on ne doit de parler » 多出介词 de，应作 « on ne doit parler »。} C’est cela qui explique la difficulté d’engager la conversation avec un inconnu puisqu’on ne sait pas quel sujet aborder avec lui. Autrement dit, on craint de violer la convention d’intérêt. Heureusement, il existe des sujets passe-partout censés intéresser tout le monde et bien commodes pour briser la glace, comme parler de la pluie et du beau temps. Il est des privilégiés qui échappent à cette loi d’intérêt : ce sont les dépositaires de l’autorité, dont la parole s’impose à leur auditoire : je parle bien sûr des … professeurs. Les mines penaudes et les regards désespérés de certains de mes étudiants pendant mon cours de linguistique me rappellent qu’il m’arrive, je le confesse, d’enfreindre cette loi.
\item Selon la loi d’informativité, un énoncé ne doit apporter à son destinataire que des informations qu’il ignore. Pourtant, en parlant de la pluie et du beau temps, on n’apprend généralement rien à son interlocuteur. Mais devant l’obligation urgente de parler pour éviter un silence gênant (dans un ascenseur, par exemple, ou face à une personne inconnue qui nous est présentée), on donne parfois la priorité à la loi d’intérêt sur la loi d’informativité.
\item La loi d’exhaustivité stipule que le locuteur est tenu de donner l’information maximale compatible avec la vérité.
\end{itemize}
Les quatre lois du discours de Ducrot sont souvent rapprochées des maximes conversationnelles du philosophe britannique Herbert Paul Grice. Vous allez voir que ce n’est pas sans raison.

\minititle{\textcircled{\scriptsize 2} Paul Grice : les maximes et implicatures conversationnelles}
L’analyse conversationnelle est l’étude des mécanismes de la conversation dans le monde social. Historiquement, son fondateur est le sociologue américain Harvey Sacks (1935-1975), selon qui la conversation est un processus minutieusement organisé possédant des règles implicites qui s’appliquent dans toute interaction indépendamment du contexte.

\minititle{1. Le principe de coopération}
En linguistique interactionnelle, le philosophe britannique \textbf{Herbert Paul Grice} (1913-1988) a bâti une théorie dans laquelle la signification réside dans la communication d’un locuteur avec autrui. C’est également le point de vue d’\textbf{Oswald Ducrot} :
\begin{lecture}{}
Il y a, dans la construction du sens, deux moments : un moment strictement linguistique, où l’on attribue une valeur à la phrase, et un second moment, que j’appelle rhétorique, où cette première valeur interagit avec la situation.
\end{lecture}
Partant du principe que la compréhension se fonde sur la conversation entre plusieurs personnes qui doivent accepter les mêmes règles, la théorie pragmatique de Paul Grice repose sur deux principes essentiels :
\begin{itemize}
\item le principe de signification naturelle : pour le destinataire, comprendre un énoncé consiste à récupérer l’intention du locuteur.\ednote{« signification naturelle » 应为 « signification non naturelle »。在 Grice 的区分中，诉诸说话者意图及听者对该意图的识别，属于非自然意义（meaning\textsubscript{NN}）；自然意义的典型例子则是症状表示疾病。\ecite{62}}
\end{itemize}
\origpage{374}{360}
\begin{itemize}
\item le principe de coopération : les implicatures (inférences) que tire le destinataire de l’énoncé sont le résultat de l’hypothèse que le locuteur coopère avec lui, c’est-à-dire participe à la conversation d’une manière efficace, raisonnable et coopérative. Grice énonce ainsi ce principe :
\end{itemize}
\begin{lecture}{}
Nous pourrions ainsi formuler en première approximation un principe général qu’on s’attendra à voir respecté par tous les participants : que votre contribution conversationnelle corresponde à ce qui est exigé de vous, au stade atteint par celle-ci, par le but ou la direction acceptés de l’échange parlé dans lequel vous êtes engagé. Ce qu’on pourrait appeler principe de coopération (cooperative principle), abrégé en CP.
\end{lecture}
Partant du premier principe, Grice a essayé de répondre à la question suivante : comment le destinataire peut-il récupérer l’intention (le « vouloir-dire ») du locuteur. Son hypothèse est qu’il y parvient à travers par la voie du principe de coopération et des maximes de conversation.\ednote{« à travers par la voie » 为两个表达的叠合，应择 « par la voie » 或 « à travers »，不应连用。}

\minititle{2. Les maximes conversationnelles}
Paul Grice édicté des maximes conversationnelles auxquelles les interlocuteurs sont tenus de se conformer.\ednote{« Paul Grice édicté » 缺谓语助动词，宜作 « Paul Grice a édicté »；也可采用现在时 « édicte »，但不能只用过去分词。} Au nombre de quatre, elles contiennent des sous-maximes et recoupent en partie les lois du discours de \textbf{Ducrot} :
\begin{itemize}
\item Maxime de quantité :
\begin{itemize}
\item Fais en sorte que ta contribution apporte autant d’informations que possible.\ednote{数量准则第一项并非“信息越多越好”，而是提供当前交谈目的所要求的充分信息（as informative as is required）。应将 « que possible » 改为 « qu’il est requis pour les buts de l’échange » 一类表述；否则与下项“不要过量”相冲突。\ecite{63}}
\item Ne donne pas plus d’informations qu’on en demande.
\end{itemize}
\item Maxime de qualité ou de véridicité (sincérité) :
\begin{itemize}
\item Ne dis pas ce que tu crois être faux.
\item Ne dis pas ce pour quoi tu manques de preuve.
\end{itemize}
\item Maxime de relation :
\begin{itemize}\item Sois pertinent.\end{itemize}
\item Maxime de manière :
\begin{itemize}
\item Évite d’être obscur.
\item Evite d’être ambigu.
\item Sois bref.
\item Sois ordonné.
\end{itemize}
\end{itemize}
Partant de l’hypothèse que le locuteur coopère avec lui, le schéma de raisonnement du destinataire de l’énoncé a la forme suivante :
\begin{itemize}
\item Le locuteur énonce A dans un contexte C.
\item La proposition A semble violer une ou plusieurs maxime (s) conversationnelle (s).
\item Le destinataire suppose que le locuteur est rationnel et coopératif.
\item Le destinataire en conclut que le locuteur viole les maximes en surface seulement, mais qu’il les respecte à un niveau plus profond.
\end{itemize}
\origpage{375}{361}
\begin{itemize}
\item Le destinataire infère une proposition B qui complète ou rectifie le contenu de A en respectant les maximes apparemment enfreintes.
\item Le destinataire en infère que le locuteur visait communiquer B. On dit alors que « B est l’implicature de A dans C ».\ednote{« visait communiquer » 缺 à，应为 « visait à communiquer »。}
\end{itemize}

\minititle{3. Les violations volontaires}
Dans une conversation, on peut être amené à violer volontairement des maximes conversationnelles. Illustrons cela avec un exemple :\ednote{Grice 区分隐蔽违反（violate）、公开不遵守以产生含意（flout）及准则冲突（clash）。本节用同一 « violation » 覆盖这些情形，容易遮蔽区别；下文“信息不足以免无据”尤其是准则冲突的例子。\ecite{63}}
\begin{quote}
\textbf{A :} Maman, tu as l’heure ?\\
\textbf{B :} Papa vient de partir au travail.
\end{quote}
Si l’on considère le sens littéral, B n’a pas répondu à la question de A et semble avoir violé le principe de relation. Mais en fait, B respecte le principe de coopération. Pour le prouver, nous allons reformuler le dialogue de la manière suivante :
\begin{quote}
\textbf{A :} Es-tu en mesure de me dire l’heure de la manière conventionnelle, indiquée sur ta montre, et si tu peux le faire, je te prie de me le dire.

\textbf{B :} Je ne sais pas en ce moment l’heure exacte, mais je peux te donner une information dont tu pourras la déduire approximativement, à savoir que papa vient de partir au travail. Il part tous les jours à la même heure, 8 heures du matin.
\end{quote}
Ainsi, si A sait (ce qui est très probable) que son père quitte la maison tous les jours à 8 heures du matin, alors il pourra inférer l’information dont il a besoin. En violant volontairement la maxime de relation, B a respecté à la fois la maxime de qualité (« Ne dis pas ce pour quoi tu manques de preuve) et de quantité (« Fais en sorte que ta contribution apporte autant d’informations que possible »).\ednote{« preuve) » 之前缺闭引号，应作 « preuve »)。关于重复出现的 « autant d’informations que possible »，参见原书 p.360 对应校注；此处仍照底本保留。}

Dans l’exemple ci-dessous, le locuteur viole la maxime de quantité pour ne pas violer la maxime de qualité :
\begin{quote}
\textbf{A :} Où habite Pierre ?\\
\textbf{B :} Quelque part en Amérique du sud.
\end{quote}
De prime abord, la réponse de B ne respecte pas la maxime de quantité, parce qu’elle ne contient pas les informations requises. Toutefois, en vertu du principe de coopération, on suppose que B a donné toutes les informations dont il disposait. En donnant une information plus précise, qui aurait satisfait la maxime de quantité, B aurait pris le risque de violer la deuxième maxime de qualité, selon laquelle qui interdit de donner des informations dont on n’est pas sûr.\ednote{« selon laquelle qui interdit » 句法叠合，宜删去 « selon laquelle »，保留 « qui interdit »。}

Il existe d’autres cas où le locuteur prononce volontairement des énoncés faux, violant ainsi intentionnellement la maxime de qualité. Illustrons ces cas à travers le dialogue suivant entre un\origcont{376}{362} fils et son père :
\begin{quote}
\textbf{A :} 2/100 à ton examen de linguistique ? Quel talent ! \hfill (Antiphrase / Ironie)

\textbf{B :} oui, j’avoue que c’est un résultat un peu faible \hfill (Euphémisme)

\textbf{A :} tu es vraiment l’étudiant le plus nul du monde ! \hfill (Hyperbole)

\textbf{B :} et toi, tu es un rat de bibliothèque ! \hfill (Métaphore)
\end{quote}
Examinons à présent le dialogue suivant :
\begin{quote}
\textbf{A :} Où est Charles ?\\
\textbf{B :} Il y a un vélo jaune devant la maison d’Anne.
\end{quote}
Si A sait que Charles a un vélo jaune, alors il va inférer que B lui suggère qu’il est possible que Charles se trouve chez Anne. Nous voyons ici que A a dû faire des inférences pour comprendre ce que B a voulu lui transmettre. Les inférences faites pour conserver le principe de coopération sont appelées par Grice implicatures conversationnelles.

\minititle{4. La notion d’implicature}
Le linguiste britannique \textbf{Gerald Gazdar} (né en 1950) définit ainsi l’implicature :
\begin{lecture}{}
Une implicature est une proposition qui est impliquée par l’assertion d’une phrase dans un contexte donné bien que cette proposition ne fasse pas partie de ce qui a été effectivement dit, ni n’en soit une conséquence logique.
\end{lecture}
Grice, quant à lui, fait une différence entre ce qui est dit (contenu vériconditionnel) et ce qui est implicité (contenu non-vériconditionnel). Il distingue aussi les implicatures conventionnelles et des implicatures conversationnelles.

\minititle{1) L’implicature conventionnelle}
Les implicatures conventionnelles sont déclenchées par les constructions linguistiques elles-mêmes. Indépendantes de l’énonciation, non-vériconditionnelles, elles dépendent de l’énoncé lui-même, ou de mots présents dans l’énoncé. Grice a donné l’exemple type suivant :
\begin{quote}
Ex : He is an Englishman; he is, therefore, brave.\\
(Implicature : His being an Englishman implies that he is brave.)
\end{quote}
Ici, l’implicature conversationnelle est déclenchée par le terme \textit{therefore}. Ce qui est implicité (l’implicature), c’est que tous les anglais sont courageaux.\ednote{本段正解释“规约含意”，故 « conversationnelle » 应为 « conventionnelle »；« courageaux » 应作 « courageux »。所给英文例句传达“其勇敢可由其英国人身份推出”的联系，并未直接断言“所有英国人都勇敢”，不宜将两者完全等同。\ecite{63}}
Regardons un autre exemple :
\begin{quote}Ex : Même Jean aime la linguistique.\end{quote}

\origpage{377}{363}
Le mot « même » va déclencher les implicatures suivantes :
\begin{itemize}
\item Jean aime la linguistique
\item D’autres personnes aiment la linguistique
\item Jean est la personne la moins susceptible d’aimer la linguistique\ednote{第一项“Jean 喜欢语言学”是此句断言的核心内容，不应同后两项一起列作含意；这也与前页对“所说”及“含意”的区分冲突。后两项的尺度、备选项及地位须结合所采理论分析，不在正文静默改写。}
\end{itemize}
Parmi les déclencheurs d’implicatures conventionnelles, on peut citer les conjonctions « mais » et « donc », mais certains linguistes incluent aussi « bien que », « quoique », « puisque », et même le verbe « réussir à » (\textit{manage}).
\begin{quote}
Ex : Marie est catholique \textit{mais} sans enfant\\
(Implicature : les catholiques ont des enfants.)

Jean a \textit{encore} raté un examen\\
(Implicature : il a déjà raté plusieurs examens.)

Martin est \textit{parvenu} à me convaincre.\\
(Implicature : ce n’est pas facile de me convaincre.)\ednote{« encore » 在这里要求已有过相应失败，而不必是“已经失败多次”：仅一次既往失败也足以使该用法成立。此外，此类重复性背景常按预设处理，不应不加说明地全部等同于 Grice 式规约含意。}
\end{quote}

\minititle{2) L’implicature conversationnelle}
Objet d’une attention particulière de Grice, l’implicature conversationnelle désigne ce qui est suggéré ou signifié par un locuteur, de façon implicite. Contrairement à l’implicature conventionnelle, l’implicature conversationnelle n’est pas une propriété de l’énoncé lui-même : elle est déclenchée par la présomption par l’interlocuteur du respect du principe de coopération et l’utilisation (ou la violation) d’une maxime de conversation.

Alors que l’implicature conventionnelle relève de la sémantique, l’implicature conversationnelle relève de la pragmatique. Si l’interlocuteur infère une proposition quelconque, on dira qu’il a tiré de l’énoncé du locuteur une implicature conversationnelle.\ednote{“推得任何命题”过宽：逻辑蕴涵、听者偶然联想到的命题等并不因此成为会话含意；还须有可归属于说话者交际用意的推理联系。Grice 的推导条件包括共享知识、合作假设及说话者可预期听者作出该推断。\ecite{63}} Grice distingue deux types d’implicatures conversationnelles :
\begin{itemize}
\item Les implicatures \textit{généralisées}, qui se constituent sans avoir besoin d’un contexte particulier :
\begin{quote}
Ex : J’ai lu quelques-uns des livres de Grice. (Implicature : je n’ai pas lu tous ses livres.)

Ce compositeur a du talent (Implicature : ce compositeur n’a pas de génie.)
\end{quote}
\item Les implicatures \textit{particularisées}, qui sont déclenchées grâce à un contexte particulier et à des connaissances particulières sur le contexte d’énonciation :
\begin{quote}
\textbf{A :} Jean n’a pas de petite amie en ce moment ?\\
\textbf{B :} Il va très souvent à Londres ces temps-ci.
\end{quote}
\end{itemize}

\origpage{378}{364}
\begin{quote}
\textbf{A :} Je suis en panne.\\
\textbf{B :} Il y a un garage au coin de la rue.
\end{quote}
On ne peut pas présenter les maximes conversationnelles de Grice sans présenter les maximes de politesse du linguiste anglais Geoffrey Leech.

\minititle{\textcircled{\scriptsize 3} Geoffrey Leech : la pragmatique de la politesse}
Discipline très récente née dans les années 70, l’étude de la politesse est un champ de recherche pluridisciplinaire qui ne relève pas que de la linguistique : la psychologie sociale, la sociologie et l’anthropologie s’y sont beaucoup intéressées. On considère généralement que l’analyse de la politesse en linguistique commence en 1973 avec la publication d’un article de la linguiste américaine \textbf{Robin Lakoff} (née en 1942) intitulé « \textit{The logic of politeness} ».

Dans ces travaux, Lakoff insiste sur l’importance du contexte social dans lequel un énoncé quelconque est émis, ainsi que les présupposés partagés de manière implicite par les participants à l’acte communicatif. Dans son article de 1973, Lakoff exprime l’idée centrale suivante :
\begin{lecture}{}
Il est considéré comme plus important dans la conversation d’éviter de déplaire à l’allocutaire que de viser la clarté du discours.
\end{lecture}
Une décennie plus tard, dans \textit{Pragmatics of politeness} (1983), le linguiste anglais \textbf{Geoffrey Leech} (1936-2014) a proposé sa propre interprétation des phénomènes liés à la politesse.\ednote{书名有误：本章原书 p.371 的延伸书目列作 \textit{Principles of pragmatics}（1983），应据该书目订正；不是此处的 \textit{Pragmatics of politeness}。} Ses maximes de politesse interagissent avec le principe de coopération de Grice et expliquent pourquoi parfois le locuteur ne respecte pas les maximes de conversation vues plus haut.

\minititle{1. Les maximes de politesse}
Basée sur les notions de coût et de bénéfice, la pragmatique de la politesse de Leech se présente sous la forme d’un ensemble de six maximes (et de leurs sous-maximes) :

\minititle{1) Maxime de tact}
\begin{quote}\itshape
Minimize the expression of beliefs which imply cost to other; maximize the expression of beliefs which imply benefit to other.
\end{quote}
\begin{itemize}
\item Minimisez le coût de l’autre.
\item Maximisez le bénéfice de l’autre.
\end{itemize}
\textbf{Ex :} \textit{Could I interrupt you for a second ?}

\minititle{2) Maxime de générosité}
\begin{quote}\itshape
Minimize the expression of beliefs that express or imply benefit to self; maximize the expression of beliefs that express or imply cost to self.
\end{quote}
\origpage{379}{365}
\begin{itemize}
\item Minimisez le bénéfice de soi.
\item Maximisez le coût de soi.
\end{itemize}
\textbf{Ex :} \textit{You relax and let me do the dishes.}

\minititle{3) Maxime d’approbation}
\begin{quote}\itshape
Minimize the expression of beliefs which express dispraise of other; maximize the expression of beliefs which express approval of other.
\end{quote}
\begin{itemize}
\item Minimisez le déplaisir de l’autre.
\item Maximisez le plaisir de l’autre.\ednote{两句法译将 dispraise / approval 误作“使对方不快／愉快”。应为“尽量减少对他人的贬抑／尽量增加对他人的赞许”；原书下一页表格的 « la critique de l’autre / l’éloge de l’autre » 也支持此订正。}
\end{itemize}
\textbf{Ex :} \textit{I know you’re a genius – would you know how to solve this math problem here?}

\Needspace{7\baselineskip}
\minititle{4) Maxime de modestie}
\begin{quote}\itshape
Minimize the expression of praise of self; maximize the expression of dispraise of self .
\end{quote}
\begin{itemize}
\item Minimisez le plaisir de soi.
\item Maximisez le déplaisir de soi.\ednote{同上，praise / dispraise 并非“快乐／不快乐”，此处应为减少自赞、增加自谦或自我贬抑。下一页表格的 « l’éloge de soi / l’autocritique » 与英文相合。}
\end{itemize}
\textbf{Ex :} \textit{Oh, I’m so stupid !}

\minititle{5) Maxime de l’accord}
\begin{quote}\itshape
Minimize the expression of disagreement between self and other; maximize the expression of agreement between self and other.
\end{quote}
\begin{itemize}
\item Minimisez le désaccord entre soi et l’autre.
\item Maximisez l’accord entre soi et l’autre.
\end{itemize}
\textbf{Ex :} \textit{Yes, but ma’am, I thought we resolved this already on your last visit.}

\minititle{6) Maxime de sympathie}
\begin{quote}\itshape
Minimize antipathy between self and other; maximize sympathy between the self and other.
\end{quote}
\begin{itemize}
\item Minimisez l’antipathie entre soi et l’autre.
\item Maximisez la sympathie entre soi et l’autre.
\end{itemize}
\textbf{Ex :} \textit{I am sorry to hear about your father.}

Dans \textit{Pragmatics of politeness}, Leech montre que des sociétés différentes ont leurs propres normes de politesse et établissent un ordre de priorité entre les maximes. Ainsi, par exemple, dans la société chinoise, les maximes de modestie et d’accord sont plus importantes que dans la société\origcont{380}{366} française.

\minititle{2. Politesse positive et politesse négative}
Leech, constatant que certains actes de langage sont intrinsèquement polis tandis que d’autres sont impolis, va distinguer politesse positive et négative :
\begin{itemize}
\item La politesse négative consiste à minimiser l’effet d’un acte de langage impoli.
\item La politesse positive consiste à maximiser l’effet d’un acte de langage poli.
\end{itemize}
\begin{center}
\small
\renewcommand{\arraystretch}{1.22}
\begin{tabularx}{\linewidth}{|>{\raggedright\arraybackslash}p{.25\linewidth}|>{\raggedright\arraybackslash}X|>{\raggedright\arraybackslash}X|}
\hline
\textbf{maxime} & \textbf{politesse négative} & \textbf{politesse positive}\\ \cline{2-3}
 & \textbf{minimiser} & \textbf{maximiser}\\ \hline
\textbf{(i) tact} & les coûts pour autrui & les bénéfices pour autrui\\ \hline
\textbf{(ii) générosité} & les bénéfices pour soi & les bénéfices pour autrui\\ \hline
\textbf{(iii) approbation} & la critique de l’autre & l’éloge de l’autre\\ \hline
\textbf{(iv) modestie} & l’éloge de soi & l’autocritique\\ \hline
\textbf{(v) accord} & le désaccord avec l’autre & l’accord avec l’autre\\ \hline
\textbf{(vi) sympathie} & l’antipathie & la sympathie\\ \hline
\end{tabularx}
\end{center}
Table: Les maximes de politesse de G. N. Leech (\textit{apud} C. Alberdi Urquizi 2009, 28).\ednote{表中 générosité 行最后一格重复了 tact 行的 « les bénéfices pour autrui »，与前页英文及法文“最大化自己的付出”不符，应为 « les coûts pour soi »。}

\minititle{3. La politesse paradoxale}
Leech mentionne également le phénomène de politesse paradoxale (\textit{politeness paradox}), qui indique un ensemble d’échanges polis pouvant conduire à l’inactivité dans une situation idéalement polie. Ainsi, deux personnes qui s’apprêtent à franchir la même porte, devraient rester immobiles, l’un cédant poliment la priorité, et l’autre refusant poliment la proposition de son partenaire.

Puisque nous sommes en Chine, prenons un exemple local : dans un restaurant, nous avons tous vu deux personnes (amis, collègues, partenaires commerciaux, etc.) se disputant (ou faisant semblant de se disputer) au moment de payer l’addition, chacun insistant pour payer et refusant de se laisser inviter par l’autre, et cela pour ne pas « perdre la face ». Cette notion de « face » nous amène justement au point suivant : les « \textit{face threatening acts} ».

\minititle{\textcircled{\scriptsize 4} Brown et Levinson : les \textit{Face Threatening Acts}}
Les travaux récents des anthropologues et linguistes \textbf{Penelope Brown} (1944-) et \textbf{Steven Levinson} (1947-) se basent sur les notions de « face » et de « territoire » proposées initialement par le sociologue et linguiste américain Erving Goffman.\ednote{« Steven Levinson » 中名字应作 Stephen；本章原书 p.371 的书目也使用 Stephen。后文同类拼写照底本保留，不再逐处重复注释。} Dans \textit{Les rites d’interaction} (1967), il définit ainsi la face :
\begin{lecture}{}
Valeur sociale positive qu’une personne revendique effectivement à travers la ligne d’action que les autres supposent qu’elle a adoptée au cours d’un contact particulier.
\end{lecture}

\Needspace{6\baselineskip}
\origpage{381}{367}
\begin{samepage}
D’après Goffman, tout contact avec d’autrui est source de conflit potentiel. Des contraintes rituelles (les règles de politesse) sont donc mises en œuvre pour protéger la face des interlocuteurs.\ednote{« avec d’autrui » 中 d’ 多余，应作 « avec autrui »。}
\par\end{samepage}

\minititle{1. Face positive et face négative}
Pour Penelope Brown et Steven Levinson, la face est constituée de deux aspects complémentaires et étroitement liés entre eux :
\begin{itemize}
\item la face négative, qui comprend les « possessions territoriales » au sens plus large du terme : territoire corporel, spatial ou temporel, biens matériels ou savoirs, secrets, etc.
\item la face positive, c’est-à-dire le narcissisme de l’individu : l’ensemble des images valorisantes que les interactants construisent et tentent d’imposer d’eux-mêmes dans l’interaction.
\end{itemize}

\minititle{2. Les menaces faites à la face}
Reprenant l’idée de Goffman, Brown et Levinson affirment eux aussi que toute rencontre humaine présente un caractère intrinsèquement menaçant pour les faces : c’est ce qu’ils appellent les « \textit{Face Threatening Acts} ». Ces FTA peuvent menacer les deux faces du locuteur ou de son interlocuteur. Dans toute interaction à deux participants, ce sont donc quatre faces qui sont mises en présence : Brown et Levinson répartissent donc les actes de langage en quatre catégories :
\begin{itemize}
\item Les actes menaçants pour la face négative de celui qui les accomplit (offres, promesses, etc.)
\item Les actes menaçants pour la face positive de celui qui les accomplit (aveux, autocritiques, etc.)
\item Les actes menaçants pour la face négative de celui qui les subit (questions personnelles, etc.)
\item Les actes menaçants pour la face positive de celui qui les subit (critiques, réfutations, etc.)
\end{itemize}

\minititle{3. Les menaces faites au territoire}
Concernant le territoire, on peut distinguer :
\begin{itemize}
\item Les actes menaçants pour le territoire de celui qui le produit : par exemple, donner un objet qui nous appartient, c’est se déposséder d’une partie de son territoire.
\item Les actes menaçants pour le territoire de celui qui reçoit : une visite à l’improviste est un exemple d’invasion du territoire d’autrui.
\end{itemize}

\minititle{4. Les notions de \textit{face want} et \textit{face work}}
Tout individu est guidé par le désir de préservation de chacune de ses deux faces : c’est ce que Brown et Levinson appellent « \textit{face want} » : d’un côté il essaie de défendre son territoire personnel et de l’autre il veut être reconnu et apprécié par les autres. Le moyen permettant de résoudre la contradiction inhérente à la volonté des interlocuteurs d’une part de s’auto-préserver, et d’autre part d’éviter de heurter les deux faces d’autrui est appelé par Goffman « \textit{face work} », terme traduit en français par « travail de figuration ». D’après Brown et Levinson, les participants y parviennent\origcont{382}{368} en mettant en œuvre diverses stratégies de politesse. Le modèle de Brown et Levinson constitue à l’heure actuelle le cadre théorique le plus cohérent en matière de politesse linguistique. Très récemment, les travaux de la linguiste française Kerbrat-Orecchioni lui ont apporté des améliorations.

\minititle{\textcircled{\scriptsize 5} Kerbrat-Orecchioni : les \textit{Face Flattering Acts}}
Pour corriger la portée excessivement négative de la théorie de Brown et Levinson, qui n’envisagent que des actes menaçants pour la face, \textbf{Catherine Kerbrat-Orecchioni} (1943-) a introduit la notion d’anti-FTA, équivalent positif du FTA, c’est-à-dire un acte de langage qui permet de préserver les faces du locuteur ou de son interlocuteur.

\minititle{1. FFA et anti FTA}
Kerbrat-Orecchioni a donné à ces actes de langage le nom de « \textit{Face Flattering Acts} » :
\begin{lecture}{}
La politesse ne se réduit pas à l’adoucissement des actes menaçants : elle peut consister, plus positivement, en la production d’actes anti-menaçants, comme les vœux ou les compliments. Ces actes valorisants pour les faces, que n’envisagent pas Brown et Levinson, nous proposons de les appeler Face Flattering Acts (ou FFA) - l’ensemble des actes de langage se répartissent alors en deux grandes familles, selon qu’ils ont sur les faces des effets essentiellement négatifs ou au contraire positifs.
\end{lecture}
Kerbrat-Orecchioni s’est également intéressée aux notions de notions de politesse négative et positive, qui sont, selon elle, sans rapport avec la face positive et négative :\ednote{« aux notions de notions de politesse » 中 « notions de » 重复，应删一组。}
\begin{itemize}
\item La politesse négative peut être « abstentionniste » : ce sont les rites d’évitement dont parle Goffman, qui consistent à ne pas commettre le FTA programmé. Elle peut aussi être « compensatoire » ou « redressive », c’est-à-dire réparatrice quand un FTA a été commis.
\item La politesse positive est « productionniste » d’anti FTA : un compliment, un simple bonjour constituent un acte de politesse positive.
\end{itemize}

\minititle{2. Les adoucisseurs de FTA}
Le travail de figuration (\textit{face work}) mentionné plus haut consiste d’une part à éviter de produire des FTA et d’autre part à produire des FFA. Si un FTA a été produit, on doit avoir recours à la politesse compensatoire et l’atténuer au moyen de \textit{softeners} (« adoucisseurs »). Parmi eux, citons l’intonation (une voix douce, un débit pas trop rapide), une mimique (un sourire), les modalisateurs (« il me semble que »), le conditionnel (« pourrais-tu fermer la porte »), la formulation indirecte (« on pourrait peut-être fermer la porte »), les excuses, l’utilisation de minimisations (« j’ai un petit service à te demander »), etc. Quant aux intensifieurs de FFA, ils sont moins diversifiés mais fonctionnent de la même façon.

Par exemple, on intensifie un remerciement ou un compliment en ajoutant certains adverbes :\origcont{383}{369} « merci \textit{beaucoup} », « ta nouvelle coiffure \textit{très} jolie », « tu es \textit{vraiment} sympa », etc. L’intonation et les mimiques peuvent aussi jouer le rôle d’intensifieur du FFA.\ednote{« ta nouvelle coiffure très jolie » 作为完整断言缺系词，宜作 « ta nouvelle coiffure est très jolie »；若为口语省略形式则需交代，本处照底本保留。}

En guise de conclusion, je vais moi aussi produire un FFA très sincère : vous êtes \textit{vraiment très} courageux, car vous venez de terminer la lecture du douzième chapitre, et donc la deuxième partie de notre livre. Vous en avez donc fini avec les six domaines des sciences du langage.

La troisième et dernière partie de notre livre (et donc les quatre chapitres à venir) sera consacrée entièrement aux différents courants de la linguistique américaine moderne et contemporaine.

\Needspace{10\baselineskip}
\phantomsection\addcontentsline{toc}{section}{Exercices}
\section*{Exercices}
\markright{Exercices}
\minititle{I. Complétez avec une interjection :}
\begin{keybox}
\small
\renewcommand{\arraystretch}{1.35}
\begin{tabularx}{\linewidth}{@{}>{\raggedright\arraybackslash}X>{\raggedright\arraybackslash}p{.30\linewidth}@{}}
\rule{20mm}{.35pt}, tu ne devrais pas être en cours, toi ? & (interpeller quelqu’un)\\
Mais \rule{20mm}{.35pt} ! Comment n’y ai-je pas pensé plus tôt ? & (exprimer une évidence)\\
\rule{20mm}{.35pt} ! Pour une surprise c’est une surprise ! & (exprimer l’étonnement)\\
\rule{20mm}{.35pt} ! Ça fait deux heures que j’essaie de les appeler ! & (exprimer la colère)\\
Le magasin est fermé, je reviendrai demain, \rule{20mm}{.35pt} ! & (exprimer la résignation)
\end{tabularx}
\end{keybox}

\minititle{II. À l’aide des exemples ci-dessous, cités par Paul Grice dans \textit{Logique et conversation}, retrouvez ses quatre maximes conversationnelles :}
\begin{keybox}
\begin{enumerate}[label=\arabic*.,leftmargin=6mm]
\item (\hspace{15mm})

J’attends une aide véritable, pas un semblant d’aide. S’il me faut du sucre pour un gâteau que quelqu’un m’aide à faire, j’espère bien qu’il ne me tendra pas le sel ; s’il me faut une cuiller, je veux croire que ce ne sera pas une attrape en caoutchouc.
\item (\hspace{15mm})

Je compte sur une aide de mon associé ajustée aux besoins immédiats de chaque stade de la transaction ; si je mélange des ingrédients pour faire un gâteau, je ne m’attends pas qu’on me tende un bon livre, ni même une pelle à tarte (même si peut-être cette contribution peut devenir opportune à un stade ultérieur).
\item (\hspace{15mm})

Si quelqu’un m’aide à réparer une voiture, je m’attends que sa contribution ne corresponde ni plus ni moins qu’à ce qui est demandé ; si par exemple à un moment donné, il me faut quatre vis, j’attends\origcont{384}{370} de lui qu’il m’en donne quatre, et non pas six ou deux.
\item (\hspace{15mm})

Je compte que mon partenaire élucide pour moi la nature de sa contribution et qu’il l’accomplisse en un temps raisonnable.
\end{enumerate}
\end{keybox}

\Needspace{14\baselineskip}
\minititle{III. Dans les dialogues suivants, dites quelles maximes de Ducros, Grice et Leech ont été utilisées ou violées :\ednote{题干 « Ducros » 为 « Ducrot » 之误；本章相应小节及书目均为 Oswald Ducrot。}}
\begin{keybox}
\renewcommand{\thempfootnote}{校\arabic{editorialnote}}
Dialogue 1 : à l’aéroport

一位贵宾对负责接待的员工

A : 王先生，这几个月你跑前跑后，真太辛苦了

B : 一点儿都不辛苦，您从国外那么远来到中国，那才是辛苦。只要您吃得好，住得舒服，那就是我最大的荣幸。

\medskip
Dialogue 2 : une rencontre entre deux inconnus

A : 请问您尊姓大名？

B : 不敢，鄙人姓柳，名云，字子强，外号小迷糊。有空请到寒舍一叙。

\medskip
Dialogue 3 : dans une salle de réunion

A : 你所提的条件，有一些，我是同意的，只是对其中的某些条款我还需要回去和同事们研究研究，再向上级请示一下。

B : 我很高兴您能表示同意，希望能早日达成协议。

\medskip
Dialogue 4 : dans le bureau du directeur

A : 主总\ednote{第四段对话称呼底本作“主总”，可能为“朱总”等姓氏误字，但现有上下文不能唯一确定，故仅列待核，不猜改人名。}，您看我的请调报告能批准吗？

B : 咱俩可是老交情了。我很理解，走，我请你喝两杯去。
\end{keybox}

\Needspace{12\baselineskip}
\minititle{IV. Trouvez les implicatures dans les énoncés suivants :}
\begin{keybox}
\begin{enumerate}[label=\arabic*.,leftmargin=6mm]
\item Je vous présente Jean, mon collègue mais néanmoins ami.
\item Marie est enceinte, mais Pierre est content.
\item Quelques étudiants ont aimé le cours de pragmatique.
\item Notre professeur de linguistique ne fait pas son âge.
\item Il est honnête, pour un banquier.
\item Même toi, tu pourrais résoudre cet exercice.
\item Pierre est riche, mais généreux.
\end{enumerate}
\end{keybox}

\Needspace{11\baselineskip}
\origpage{385}{371}
\phantomsection\addcontentsline{toc}{section}{Pour aller plus loin}
\section*{Pour aller plus loin}
\markright{Pour aller plus loin}
\begingroup\small\setstretch{1.04}
\begin{list}{}{\setlength{\leftmargin}{4mm}\setlength{\itemindent}{-4mm}\setlength{\itemsep}{5pt}}
\item BROWN (Penelope) et LEVINSON (Stephen). \textit{Politeness: Some Universals in Language Usage. Studies in interactional sociolinguistics.} Cambridge University Press, Cambridge, 1987.
\item DUCROT (Oswald). \textit{Le Dire et le dit.} Éditions de Minuit, Paris, 1984.
\item GAZDAR (Gerald James Michael). \textit{Formal pragmatics for natural language implicature, presupposition and logical form.} Ph.D. diss.. University of Reading. Reading (Royaume-Uni), 1976.
\item GRICE (H. Paul). \textit{Meaning.} The Philosophical Review, Duke University Press, Durham (Caroline du Nord, États-Unis), 1957, 66(3):377–388.
\item GRICE (H. Paul). \textit{Logic and conversation.} In Peter Cole, \& Jerry L. Morgan (Eds.), Syntax and Semantics, Vol. 3, Speech Acts, Academic Press, New York, 1975, 41-58p.
\item HALLIDAY (Michael Alexander Kirkwood) et HASAN (Ruqaiya). \textit{Cohesion in English.} Longman, Londres, 1976.
\item KERBRAT-ORECCHIONI (Catherine). \textit{Le discours en interaction.} A. Colin, Paris, 2005.
\item KERBRAT-ORECCHIONI (Catherine). \textit{Les interactions verbales III,} A. Colin, Paris, 1998.
\item KESIK (Marek). \textit{La cataphore.} Presses Universitaires de France, Paris, 1989.
\item LAKOFF (Robin). \textit{The logic of politeness: Or, minding your p’s and q’s.} In C. Corum, T. Cedric Smith-Stark, \& A. Weiser (Eds.), Papers from the \textit{9th Regional Meeting of the Chicago Linguistic Society,} Chicago Linguistic Society, Chicago, 1973.
\item LEECH (Geoffrey). \textit{Principles of pragmatics.} London ; New York: Longman Group Ltd, 1983.
\item LEVINSON (Stephen C.). \textit{Pragmatics.} Cambridge University Press, Cambridge, 1983
\item MAILLARD (Michel). « \textit{Essai de typologie des substituts diaphoriques} », in \textit{Langue Française,} 1974.
\end{list}
\endgroup

\needspace{65mm}
\phantomsection\addcontentsline{toc}{section}{Glossaire}
\section*{Glossaire}
\markright{Glossaire}
\origpage{386}{372}
\gloss{énoncé 陈述}{描述事态、行动、情感或信念的语句。}
\gloss{déictique 指示成分}{词或短语将话语直接与时间、地点或人物联系起来的一种说法。}
\gloss{performativité 行事话语}{言语行为理论中执行动作的语句。\ednote{法文名词 performativité 指“施为性／行事性”，不是执行行为的“语句”本身。若保留此中文定义，法文词头宜为 « énoncé performatif »；若保留词头，则定义应改为话语实施行为的性质。参见本章前文对 énoncé performatif 的说明。}}
\gloss{situation d’énonciation 情景语境}{一个单词、话语或者文本出现的语言语境和情景语境。话语的意思不但要通过单词的字面意思也要靠它们出现的语境或情景来确定。}
